\documentclass[
  10pt,
  letterpaper,
  twoside,
  openright,
  numbers=noenddot,
  headings=standardclasses
]{scrbook}

\usepackage{geometry}
\geometry{
  paperwidth=8in,
  paperheight=10in,
  inner=0.875in,
  outer=1.75in,
  top=0.875in,
  bottom=0.875in,
  marginparwidth=1.25in,
  marginparsep=0.25in,
  headheight=14pt,
  headsep=0.25in,
  footskip=0.4in
}

\usepackage{fontspec}
\usepackage{xcolor}
\usepackage{graphicx}
\usepackage{tikz}
\usetikzlibrary{calc, positioning}
\usepackage{amsmath, amssymb}
\usepackage{booktabs}
\usepackage{microtype}
\usepackage{tcolorbox}
\tcbuselibrary{skins, breakable}
\usepackage{hyperref}
\usepackage{bookmark}
\usepackage{sidenotes}
\usepackage{caption}

% =============================================================================
% PHYSICAL AI THEME COLORS & PALETTE
% =============================================================================
% Sets the primary accent color and semantic layer for Physical AI
% Include this file BEFORE header-includes.tex
% =============================================================================

% Primary branding colors
\definecolor{harvardcrimson}{HTML}{A51C30}
\definecolor{ethdarkblue}{HTML}{1F407A}
\definecolor{ethblue}{HTML}{215CAF}
\definecolor{ethpetrol}{HTML}{007A87}
\definecolor{ethbronze}{HTML}{B87333}
\definecolor{ethpurple}{HTML}{5B4B8A}

% Canonical MLSysBook variables
\definecolor{accentcolor}{HTML}{1A4D3E}     % Primary accent for headings and crossrefs (Emerald Pine - Complement to Crimson)
\definecolor{toccolor}{HTML}{000000}        % TOC entries: black (matches Vol1/Vol2)

% Semantic layer for Physical AI diagrams, links, and callouts
\colorlet{structural}{ethdarkblue}
\colorlet{measured}{ethpetrol}
\colorlet{proposal}{ethbronze}
\colorlet{consequence}{harvardcrimson}
\colorlet{authority}{ethpurple}

% Ink levels for body, secondary, and hairline material
\definecolor{ink}{HTML}{1A202C}
\definecolor{softink}{HTML}{4A5568}
\definecolor{paleink}{HTML}{A0AEC0}

% Surfaces
\definecolor{cardbg}{HTML}{F8FAFC}
\definecolor{shadecolor}{HTML}{F4F6F9}

% Quarto standard callouts harmonized to the palette
\definecolor{quarto-callout-note-color}{HTML}{1F407A}
\definecolor{quarto-callout-note-color-frame}{HTML}{1F407A}
\definecolor{quarto-callout-tip-color}{HTML}{007A87}
\definecolor{quarto-callout-tip-color-frame}{HTML}{007A87}
\definecolor{quarto-callout-warning-color}{HTML}{A51C30}
\definecolor{quarto-callout-warning-color-frame}{HTML}{A51C30}
\definecolor{quarto-callout-important-color}{HTML}{5B4B8A}
\definecolor{quarto-callout-important-color-frame}{HTML}{5B4B8A}
\definecolor{quarto-callout-caution-color}{HTML}{A51C30}
\definecolor{quarto-callout-caution-color-frame}{HTML}{A51C30}

% ===========================================================================
% Physical AI Systems, LaTeX preamble
% ===========================================================================
%
% Engine:  LuaLaTeX
% Class:   KOMA-Script scrbook (set in _quarto.yml)
% Loaded:  first of the include-in-header files. tex/cover.tex follows and
%          depends on the palette defined here.
%
% RULES FOR EDITING THIS FILE
%
%   1. Every setting is declared exactly once. To change a value, change it
%      where it lives. Do not append an override at the bottom. This file
%      previously carried two \setlist declarations, two \arraystretch values,
%      two \emergencystretch values, and fvextra loaded twice, all because
%      changes were bolted on rather than integrated.
%   2. Sections run in dependency order: packages, palette, page, type,
%      floats, code, notes, links, utilities. Put new settings in the section
%      they belong to.
%   3. Anything that draws something belongs in cover.tex or in a figure
%      source, not here.
%   4. Comment the intent, not the syntax. A future reader needs to know why a
%      value is what it is, not that \setlength sets a length.
%
% Page geometry, font size, and leading live in _quarto.yml, because Quarto
% owns the geometry package options. This file assumes them, never sets them.
%
% ===========================================================================


% ---------------------------------------------------------------------------
% 1. Packages
% ---------------------------------------------------------------------------

\usepackage{microtype}          % Optical margins and glyph expansion.
\usepackage{float}              % Strict float placement [H].
\usepackage{morefloats}         % Expand LaTeX float registers.
\usepackage{placeins}           % Prevent floats from overshooting section/chapter boundaries.
\usepackage{needspace}          % Dynamic space reservation to prevent awkward table breaks.
\usepackage{scrlayer-scrpage}   % KOMA page styles, for running heads.
\usepackage{etoolbox}           % Patching KOMA and Quarto internals.
\usepackage{array}              % Column specifiers used by pandoc tables.
\usepackage{booktabs}           % \toprule and friends in pandoc tables.
\usepackage{enumitem}           % List spacing control.
\usepackage{framed}             % Backing for the Shaded code environment.
\usepackage{fancyvrb}           % Verbatim engine that pandoc emits into.
\usepackage{fvextra}            % Line breaking inside listings.
\usepackage{ragged2e}            % Balanced ragged-right typesetting with hyphenation.
\usepackage{marginnote}         % Margin column content.
\usepackage{environ}            % \NewEnviron, used by the fullwidth utility.
\usepackage{lettrine}           % Drop caps for chapter openers.
\usepackage{colortbl}           % Alternating row tints and table styling.
\usepackage{tikz}               % Cover artwork.
\usetikzlibrary{arrows.meta,positioning,shapes.geometric,fit,backgrounds,calc}


% ---------------------------------------------------------------------------
% 2. Palette
%
% Harvard crimson and ETH Zurich blue, plus a semantic layer. Prose, callouts,
% and figures reference the semantic names, so a palette change happens in one
% place. The raw names stay because the TikZ figure sources use them directly.
% ---------------------------------------------------------------------------

\definecolor{harvardcrimson}{HTML}{A51C30}
\definecolor{ethdarkblue}{HTML}{1F407A}
\definecolor{ethblue}{HTML}{215CAF}
\definecolor{ethpetrol}{HTML}{007A87}
\definecolor{ethbronze}{HTML}{B87333}
\definecolor{ethpurple}{HTML}{5B4B8A}

% Semantic layer. These encode the book's argument: what is structural, what
% is measured, what is only proposed, what has consequences, who has authority.
\colorlet{structural}{ethdarkblue}
\colorlet{measured}{ethpetrol}
\colorlet{proposal}{ethbronze}
\colorlet{consequence}{harvardcrimson}
\colorlet{authority}{ethpurple}

% Ink levels for body, secondary, and hairline material.
\definecolor{ink}{HTML}{1A202C}
\definecolor{softink}{HTML}{4A5568}
\definecolor{paleink}{HTML}{A0AEC0}

% Surfaces. shadecolor is named by the framed/Shaded contract, not by us.
\definecolor{cardbg}{HTML}{F8FAFC}
\definecolor{shadecolor}{HTML}{F4F6F9}

% Quarto's standard callouts, harmonized to the semantic layer.
\definecolor{quarto-callout-note-color}{HTML}{1F407A}
\definecolor{quarto-callout-note-color-frame}{HTML}{1F407A}
\definecolor{quarto-callout-tip-color}{HTML}{007A87}
\definecolor{quarto-callout-tip-color-frame}{HTML}{007A87}
\definecolor{quarto-callout-warning-color}{HTML}{A51C30}
\definecolor{quarto-callout-warning-color-frame}{HTML}{A51C30}
\definecolor{quarto-callout-important-color}{HTML}{5B4B8A}
\definecolor{quarto-callout-important-color-frame}{HTML}{5B4B8A}
\definecolor{quarto-callout-caution-color}{HTML}{A51C30}
\definecolor{quarto-callout-caution-color-frame}{HTML}{A51C30}


% ---------------------------------------------------------------------------
% 3. Page furniture (Running headers & footers)
%
% Standard two-sided academic textbook layout (Cambridge / MIT Press):
%   Even (verso / left) pages:  [Page #] ........................ [Chapter Title]
%   Odd (recto / right) pages:  [Section Title] ........................ [Page #]
% Subtitles and linebreaks are strictly stripped from headers.
% ---------------------------------------------------------------------------

\clearpairofpagestyles

\KOMAoptions{headwidth=textwithmarginpar,footwidth=textwithmarginpar}
\KOMAoptions{headsepline=0.4pt}
\KOMAoptions{markcase=noupper}
\addtokomafont{headsepline}{\color{accentcolor}}

\renewcommand*{\chaptermarkformat}{}
\renewcommand*{\sectionmarkformat}{}

\setkomafont{pagehead}{\normalfont\rmfamily\fontsize{8.5pt}{10.5pt}\selectfont\color{accentcolor}}
\setkomafont{pagenumber}{\normalfont\rmfamily\fontsize{8.5pt}{10.5pt}\selectfont\color{accentcolor}}

% Verso (even / left) header: Left is Folio, Right is Section mark (matching MLSysBook)
\lehead{\pagemark}
\rehead{\small\rightmark}

% Recto (odd / right) header: Left is Chapter mark, Right is Folio (matching MLSysBook)
\lohead{\small\leftmark}
\rohead{\pagemark}

% Chapter opening pages (plain style): centered folio at bottom in accent color
\cfoot[\pagemark]{}
\setkomafont{pagefoot}{\normalfont\rmfamily\fontsize{8.5pt}{10.5pt}\selectfont\color{accentcolor}}

\automark[section]{chapter}



% ---------------------------------------------------------------------------
% 4. Typography
%
% Serif body at a 4.85in measure, with matching classical serif headings
% matching MLSysBook Volumes 1 and 2.
%
% Paragraphs are indented rather than separated by vertical space. Vertical
% space marks structure here, not every paragraph break.
% ---------------------------------------------------------------------------

\KOMAoptions{fontsize=9.5pt}

\frenchspacing
\microtypesetup{protrusion=true, expansion=true, final}

\setlength{\parskip}{0pt plus 0.6pt}
\setlength{\parindent}{1.2em}

% Justification. Long technical compounds at this measure produce ladders of
% hyphens. Loosen the badness tolerance rather than the word spacing, so the
% grey of the page stays even.
\hyphenpenalty=180
\tolerance=1400
\setlength{\emergencystretch}{2.2em}

% Widow and orphan control (Cambridge / MIT Press standard)
\clubpenalty=10000
\widowpenalty=10000
\displaywidowpenalty=10000
\raggedbottom

% Pandoc emits \linebreak inside margin material where it is invalid.
\let\origlinebreak\linebreak
\renewcommand{\linebreak}[1][4]{\ifhmode\origlinebreak[#1]\fi}

% --- Headings (Matching MLSysBook Volumes 1 and 2) ------------------------

\setkomafont{disposition}{\normalfont\rmfamily\color{accentcolor}}
\setkomafont{chapter}{\normalfont\rmfamily\bfseries\fontsize{20pt}{24pt}\selectfont\color{accentcolor}}
\setkomafont{section}{\normalfont\rmfamily\bfseries\fontsize{13pt}{16pt}\selectfont\color{accentcolor}}
\setkomafont{subsection}{\normalfont\rmfamily\bfseries\fontsize{11pt}{14pt}\selectfont\color{accentcolor}}
\setkomafont{subsubsection}{\normalfont\rmfamily\bfseries\fontsize{10pt}{13pt}\selectfont\color{accentcolor}}
\setkomafont{paragraph}{\normalfont\rmfamily\bfseries\fontsize{9.5pt}{12pt}\selectfont\color{accentcolor}}

% Drop caps configuration for chapter openings (Matching MLSysBook)
\renewcommand{\LettrineFontHook}{\rmfamily\bfseries\color{accentcolor}}
\renewcommand{\LettrineTextFont}{\normalfont\scshape}
\setlength{\DefaultFindent}{2pt}
\setlength{\DefaultNindent}{0pt}

% Display math spacing
\setlength{\abovedisplayskip}{6pt plus 2pt minus 2pt}
\setlength{\belowdisplayskip}{6pt plus 2pt minus 2pt}
\setlength{\abovedisplayshortskip}{0pt plus 2pt}
\setlength{\belowdisplayshortskip}{4pt plus 2pt minus 2pt}

% Chapter subtitle command for chapter openers (appears before the horizontal rule)
\DeclareRobustCommand{\chaptersubtitle}[1]{%
  \texorpdfstring{%
    \newline{\normalfont\fontsize{11pt}{14pt}\selectfont\color{ethdarkblue}\textit{#1}}%
  }{}%
}

% --- Chapter & Section Formatting (Matching Volume 1 & Volume 2) ---------
\usepackage[explicit]{titlesec}

% Giant chapter-number font in margin
\newcommand{\mlsysChapNumFont}{\normalfont\rmfamily\fontsize{70pt}{70pt}\selectfont\bfseries}

% Numbered chapters (Large serif title with giant margin number)
\titleformat{\chapter}[block]
  {\normalfont\rmfamily\bfseries\color{accentcolor}}
  {}
  {0pt}
  {%
    \begin{tikzpicture}[remember picture, overlay]
      \node[anchor=north east,
            inner sep=0pt,
            font={\mlsysChapNumFont},
            text=accentcolor]
        at ([xshift=-40pt]current page.north east |- current page text area.north)
        {\thechapter};
    \end{tikzpicture}%
    {\fontsize{34pt}{38pt}\selectfont #1}%
  }
  []

% Unnumbered chapters (Preface, Notation, Syllabus): clean large serif title without horizontal bars
\titleformat{name=\chapter,numberless}
  {\normalfont\rmfamily\huge\bfseries\color{accentcolor}}
  {}
  {0pt}
  {\Huge #1}
  []

\titlespacing*{\chapter}{0pt}{0pt}{24pt}

% Section (Large, bold serif, accentcolor)
\titleformat{\section}
  {\normalfont\rmfamily\Large\bfseries\color{accentcolor}\raggedright}
  {\thesection}
  {0.5em}
  {#1}
\titlespacing*{\section}{0pc}{12pt plus 3pt minus 3pt}{4pt plus 1pt minus 1pt}[0pc]

% Subsection (large, bold serif, accentcolor)
\titleformat{\subsection}
  {\normalfont\rmfamily\large\bfseries\color{accentcolor}\raggedright}
  {\thesubsection}
  {0.5em}
  {#1}
\titlespacing*{\subsection}{0pc}{10pt plus 3pt minus 3pt}{4pt plus 1pt minus 1pt}[0pc]

% Subsubsection (normal size, bold serif, accentcolor)
\titleformat{\subsubsection}
  {\normalfont\rmfamily\normalsize\bfseries\color{accentcolor}\raggedright}
  {\thesubsubsection}
  {0.5em}
  {#1}
\titlespacing*{\subsubsection}{0pc}{9pt plus 3pt minus 3pt}{3pt plus 1pt minus 1pt}[0pc]

% --- Lists ----------------------------------------------------------------

\setlist[itemize]{leftmargin=1.1em,topsep=2pt,itemsep=1pt,parsep=0pt}
\setlist[enumerate]{leftmargin=1.3em,topsep=2pt,itemsep=1pt,parsep=0pt}

% --- Table of Contents & Structure Dividers (MLSysBook Modern Layout) ------

\makeatletter
% State flags for TOC divisions
\newif\ifmlsysinbackmatter
\mlsysinbackmatterfalse
\newif\ifsynthesiswritten
\synthesiswrittenfalse
\newif\if@firstnumbered
\@firstnumberedtrue
\newcounter{mlsysappendixcount}

% Pure TeX helper to distinguish Chapter numbers (1..17) from Appendix letters (A..G)
\newcommand{\format@toc@number}[1]{%
  \if\relax\detokenize\expandafter{\romannumeral-0#1}\relax
    Chapter~#1%
  \else
    Appendix~#1%
  \fi
}

% Start frontmatter with Roman numerals and suppress subsection noise in frontmatter
\AtBeginDocument{%
  \pagenumbering{roman}%
  \renewcommand{\thepage}{\roman{page}}%
  \renewcommand*\contentsname{{\normalfont\bfseries\fontsize{18pt}{22pt}\selectfont\color{consequence}Contents}}%
  \addtocontents{toc}{\protect\setcounter{tocdepth}{0}}%
  % Hook addcontentsline AFTER hyperref initialization
  \let\orig@addcontentsline\addcontentsline
  \renewcommand{\addcontentsline}[3]{%
    \Ifstr{#1}{toc}{%
      \Ifstr{#2}{part}{%
        \ifmlsysinbackmatter
          % Suppress Appendices part divider; Back Matter division handles it
        \else
          \orig@addcontentsline{#1}{#2}{#3}%
        \fi
      }{%
        \Ifstr{#2}{chapter}{%
          \ifmlsysinbackmatter
            \Ifstr{#3}{References}{%
              \orig@addcontentsline{#1}{#2}{#3}%
            }{%
              \Ifstr{#3}{Index}{%
                \orig@addcontentsline{#1}{#2}{#3}%
              }{%
                \stepcounter{mlsysappendixcount}%
                \orig@addcontentsline{toc}{chapter}{\protect\numberline{\Alph{mlsysappendixcount}}#3}%
              }%
            }%
          \else
            \orig@addcontentsline{#1}{#2}{#3}%
          \fi
        }{%
          \orig@addcontentsline{#1}{#2}{#3}%
        }%
      }%
    }{%
      \orig@addcontentsline{#1}{#2}{#3}%
    }%
  }%
}

\let\orig@appendix\appendix
\renewcommand{\appendix}{%
  \orig@appendix
  \global\mlsysinbackmattertrue
  \addtocontents{toc}{\protect\tocdivisionentry{Back Matter}}%
  \addtocontents{toc}{\protect\setcounter{tocdepth}{0}}%
}

% Patch \mainmatter so Quarto does not prematurely switch to Arabic before Chapter 1
\let\orig@mainmatter\mainmatter
\renewcommand{\mainmatter}{%
  % Defer pagenumbering{arabic} until Chapter 1 is reached
}

\newif\ifiscontentstoc
\iscontentstocfalse

% Patch \tableofcontents to guarantee Crimson/Accent "Contents" heading
\let\orig@tableofcontents\tableofcontents
\renewcommand{\tableofcontents}{%
  \begingroup
    \iscontentstoctrue
    \DeclareRobustCommand{\chaptersubtitle}[1]{}%
    \renewcommand*\contentsname{{\normalfont\huge\bfseries\color{accentcolor}Contents}}%
    \orig@tableofcontents
  \endgroup
}

% Major division banners in TOC (e.g. "Main Text", "Synthesis", "Back Matter")
\DeclareRobustCommand{\tocdivisionentry}[1]{%
  \par\addvspace{18pt}%
  \noindent\hfil{\normalfont\sffamily\bfseries\fontsize{10.5pt}{13.5pt}\selectfont\color{ink}#1}\hfil\par%
  \vspace{3pt}%
  \noindent\hfil{\color{paleink!80!softink}\rule{0.45\textwidth}{0.45pt}}\hfil\par%
  \addvspace{8pt}%
}

% Part banners in TOC (e.g. "Part I: The Machine")
\DeclareRobustCommand{\tocpartentry}[1]{%
  \par\addvspace{16pt}%
  \noindent\hfil{\normalfont\sffamily\bfseries\fontsize{10pt}{13pt}\selectfont\color{consequence}#1}\hfil\par%
  \addvspace{7pt}%
}

% Hook into \addparttocentry to emit clean centered Crimson Part entries and division lines
\renewcommand*{\addparttocentry}[2]{%
  \Ifstr{#1}{}{%
    \Ifstr{#2}{Appendices}{}{%
      \addtocontents{toc}{\protect\tocdivisionentry{#2}}%
    }%
  }{%
    \Ifstr{#2}{}{%
      \addtocontents{toc}{\protect\tocpartentry{Part~#1}}%
    }{%
      \addtocontents{toc}{\protect\tocpartentry{#2}}%
    }%
  }%
}

% Suppress raw Appendices part entry and format part lines cleanly in TOC
\renewcommand*{\l@part}[2]{%
  \in@{Part}{#1}%
  \ifin@
    \par\addvspace{16pt}%
    \noindent\hfil{\normalfont\sffamily\bfseries\fontsize{10pt}{13pt}\selectfont\color{consequence}%
      \begingroup
        \renewcommand{\numberline}[1]{}%
        #1%
      \endgroup
    }\hfil\par%
    \addvspace{7pt}%
  \fi
}

% Intercept \chapter command to cleanly handle frontmatter vs mainmatter transitions
\let\old@chapter\chapter
\renewcommand{\chapter}{%
  \@ifstar{\unnumbered@chapter}{\numbered@chapter}%
}

\newcommand{\cleanmark}[1]{%
  \begingroup
    \def\chaptersubtitle##1{}%
    \def\newline{}%
    \def\texorpdfstring##1##2{##2}%
    \protected@xdef\@clean@title{#1}%
  \endgroup
  \markboth{\@clean@title}{}%
}

\newcommand{\unnumbered@chapter}[1]{%
  \cleardoublepage
  \old@chapter*{#1}%
  \cleanmark{#1}%
}

\newcommand{\numbered@chapter}[2][]{%
  \cleardoublepage
  \if@firstnumbered
    \pagenumbering{arabic}%
    \renewcommand{\thepage}{\arabic{page}}%
    \setcounter{page}{1}%
    \addtocontents{toc}{\protect\tocdivisionentry{Main Text}}%
    \addtocontents{toc}{\protect\setcounter{tocdepth}{1}}%
    \global\@firstnumberedfalse
  \fi
  \refstepcounter{chapter}%
  \setcounter{section}{0}%
  \if\relax\detokenize{#1}\relax
    \old@chapter*{{\chapterformat}#2}%
    \cleanmark{#2}%
    \addcontentsline{toc}{chapter}{\protect\numberline{\thechapter}#2}%
  \else
    \old@chapter*[#1]{{\chapterformat}#2}%
    \cleanmark{#1}%
    \addcontentsline{toc}{chapter}{\protect\numberline{\thechapter}#1}%
  \fi
}

% Definitive Chapter Entry layout in TOC (MLSysBook layout)
\renewcommand*{\l@chapter}[2]{%
  \ifnum \c@tocdepth >\m@ne
    \addpenalty{-\@highpenalty}%
    \vskip 0.60em\relax
    {%
      \leftskip 0pt\relax
      \rightskip \@tocrmarg
      \parfillskip -\rightskip
      \parindent 0pt\relax
      \@afterindenttrue
      \interlinepenalty\@M
      \leavevmode
      \begingroup
        \renewcommand{\numberline}[1]{{\normalfont\rmfamily\bfseries\fontsize{9.5pt}{12pt}\selectfont\color{ink}\format@toc@number{##1}\hspace{0.55em}}}%
        {\normalfont\rmfamily\bfseries\fontsize{9.5pt}{12pt}\selectfont\color{ink}#1}\nobreak
      \endgroup
      \TOCLineLeaderFill
      \nobreak
      {\normalfont\rmfamily\bfseries\fontsize{9.5pt}{12pt}\selectfont\color{ink}#2}%
      \par
    }%
  \fi
}

% Definitive Section Entry layout in TOC (MLSysBook layout)
\renewcommand*{\l@section}[2]{%
  \ifnum \c@tocdepth >\z@
    \vskip 0.10em\relax
    {%
      \leftskip 1.4em\relax
      \rightskip \@tocrmarg
      \parfillskip -\rightskip
      \parindent 0pt\relax
      \@afterindenttrue
      \interlinepenalty\@M
      \leavevmode
      \begingroup
        \renewcommand{\numberline}[1]{\makebox[2.9em][l]{##1\hspace{0.4em}}}%
        {\normalfont\rmfamily\fontsize{8.8pt}{11pt}\selectfont\color{ink}#1}\nobreak
      \endgroup
      \TOCLineLeaderFill
      \nobreak
      {\normalfont\rmfamily\fontsize{8.8pt}{11pt}\selectfont\color{ink}#2}%
      \par
    }%
  \fi
}
\makeatother


% ---------------------------------------------------------------------------
% 5. Floats: captions and tables
%
% Captions lead with a bold takeaway, so the label sits in structural blue and
% the caption text drops to secondary ink.
% ---------------------------------------------------------------------------

\setkomafont{captionlabel}{\normalfont\sffamily\bfseries\fontsize{8.2pt}{10.5pt}\selectfont\color{structural}}
\setkomafont{caption}{\normalfont\sffamily\fontsize{8.0pt}{10.2pt}\selectfont\color{softink}}

% Tables: Clean, academic, compact styling (Cambridge/MIT Press monograph style)
\renewcommand{\arraystretch}{1.18}
\setlength{\tabcolsep}{3.8pt}

% Table rules styling
\setlength{\heavyrulewidth}{1.0pt}
\setlength{\lightrulewidth}{0.5pt}

% Prevent awkward orphan headers: require at least 6 baselines of vertical space
% before starting a longtable, allowing tables to render directly under their introducing prose
\BeforeBeginEnvironment{longtable}{%
  \par\penalty0\needspace{6\baselineskip}%
}
\AtBeginEnvironment{longtable}{%
  \normalfont\sffamily\fontsize{7.6pt}{9.6pt}\selectfont%
}
\AfterEndEnvironment{longtable}{%
  \normalfont\selectfont%
}
\AtBeginEnvironment{tabular}{%
  \normalfont\sffamily\fontsize{7.6pt}{9.6pt}\selectfont%
  \setlength{\tabcolsep}{3.8pt}%
  \renewcommand{\arraystretch}{1.18}%
}
\AfterEndEnvironment{tabular}{%
  \normalfont\selectfont%
}
\AtBeginEnvironment{tabular*}{%
  \normalfont\sffamily\fontsize{7.6pt}{9.6pt}\selectfont%
  \setlength{\tabcolsep}{3.8pt}%
  \renewcommand{\arraystretch}{1.18}%
}
\AfterEndEnvironment{tabular*}{%
  \normalfont\selectfont%
}

% Floats: Academic textbook standard float placement with placeins barrier
\makeatletter
\def\fps@figure{tbp}
\def\fps@table{htbp}
\makeatother
\setcounter{topnumber}{3}
\setcounter{bottomnumber}{2}
\setcounter{totalnumber}{5}
\setcounter{dbltopnumber}{3}
\renewcommand{\topfraction}{0.90}       % max fraction of page for floats at top
\renewcommand{\bottomfraction}{0.60}    % allow bottom floats
\renewcommand{\textfraction}{0.10}      % min fraction of page for text
\renewcommand{\floatpagefraction}{0.65} % min fraction of float page occupied
\renewcommand{\dbltopfraction}{0.90}
\renewcommand{\dblfloatpagefraction}{0.65}
\setlength{\floatsep}{10pt plus 2pt minus 2pt}
\setlength{\textfloatsep}{12pt plus 2pt minus 4pt}
\setlength{\intextsep}{10pt plus 2pt minus 2pt}
\setlength{\dblfloatsep}{10pt plus 2pt minus 2pt}
\setlength{\dbltextfloatsep}{12pt plus 2pt minus 4pt}






% ---------------------------------------------------------------------------
% 6. Code listings
%
% commandchars is load bearing. Pandoc emits \CommentTok{...} and
% \NormalTok{...} inside Highlighting, and without commandchars those macros
% are typeset as literal text instead of applied. Redefining Highlighting
% without it silently breaks every syntax-highlighted block in the book,
% including the dossier YAML, which is exactly what happened before.
% ---------------------------------------------------------------------------

% Shaded environment: tighten padding and frame for listings
\ifcsname Shaded\endcsname\else
  \newenvironment{Shaded}{}{}
\fi
\renewenvironment{Shaded}{%
  \begin{tcolorbox}[
    enhanced,
    breakable,
    colback=paleink!10!white,
    colframe=paleink!30!white,
    boxrule=0.4pt,
    arc=1.0pt,
    left=3pt,
    right=3pt,
    top=2pt,
    bottom=2pt,
    before skip=3pt,
    after skip=4pt
  ]%
}{%
  \end{tcolorbox}%
}

\DefineVerbatimEnvironment{Highlighting}{Verbatim}{
  commandchars=\\\{\},
  fontsize=\fontsize{6.6pt}{8.0pt}\ttfamily\selectfont,
  breaklines=true,
  breakanywhere=true,
  xleftmargin=1pt,
  xrightmargin=1pt
}

% Plain verbatim carries no tokens, so it takes no commandchars.
\RecustomVerbatimEnvironment{verbatim}{Verbatim}{
  fontsize=\fontsize{6.6pt}{8.0pt}\ttfamily\selectfont,
  breaklines=true,
  breakanywhere=true,
  xleftmargin=1pt,
  xrightmargin=1pt
}


% ---------------------------------------------------------------------------
% 7. Margin notes and footnotes
%
% Footnotes are set in the margin (reference-location: margin in _quarto.yml),
% so both read as one register: small, sans, secondary ink, sized compactly
% for the narrow margin column.
% ---------------------------------------------------------------------------

\renewcommand*{\marginfont}{%
  \RaggedRight\normalfont\sffamily\fontsize{7.5pt}{9.8pt}\selectfont\color{softink}}
\setkomafont{footnote}{%
  \normalfont\sffamily\fontsize{7.8pt}{10.2pt}\selectfont\color{softink}}
\setkomafont{footnotemark}{%
  \normalfont\sffamily\bfseries\fontsize{7.8pt}{10.2pt}\selectfont\color{consequence}}

% Sidenotes package hook: Quarto emits \sidenote{\footnotesize ...} when
% reference-location: margin is used. We define custom handlers to neutralize
% pandoc's \footnotesize and enforce compact sans-serif typography in margin notes.
\makeatletter
\NewDocumentCommand \customsidenotetext { o o +m }
{%
  \IfNoValueOrEmptyTF{#1}%
    {%
      \marginnote{%
        \RaggedRight\normalfont\sffamily\fontsize{6.8pt}{8.8pt}\selectfont\color{softink}%
        {\bfseries\color{consequence}\textsuperscript{\thesidenote}}\nobreakspace
        {\let\footnotesize\relax #3}\par
      }%
      \refstepcounter{sidenote}%
    }%
    {%
      \marginnote{%
        \RaggedRight\normalfont\sffamily\fontsize{6.8pt}{8.8pt}\selectfont\color{softink}%
        {\bfseries\color{consequence}\textsuperscript{#1}}\nobreakspace
        {\let\footnotesize\relax #3}\par
      }%
    }%
}

\NewDocumentCommand \customsidenotemark { m }
{%
  \leavevmode
  \ifhmode
      \edef \@x@sf {\the \spacefactor }%
      \nobreak
  \fi
  \hbox {\textsuperscript{\normalfont\sffamily\bfseries\fontsize{6.8pt}{8.8pt}\selectfont\color{consequence}#1}}%
  \ifhmode
      \spacefactor \@x@sf
  \fi
  \relax
}

\AtBeginDocument{%
  \renewcommand*{\@sidenotes@placemarginal}[2]{%
    \IfNoValueOrEmptyTF{#1}{\marginnote{#2}}{\marginnote[#1]{#2}}%
  }%
  \let\sidenotetext\customsidenotetext
  \let\@sidenotes@thesidenotemark\customsidenotemark
}
\makeatother


% ---------------------------------------------------------------------------
% 8. Links
%
% Deferred to \AtBeginDocument so hyperref has loaded and Quarto has finished
% its own \hypersetup.
% ---------------------------------------------------------------------------

\AtBeginDocument{%
  \hypersetup{
    colorlinks=true,
    linkcolor=ethblue,
    urlcolor=ethpetrol,
    citecolor=ethpetrol,
    filecolor=ethpetrol
  }%
}


% ---------------------------------------------------------------------------
% 9. Utilities & Margin Definition Badges
% ---------------------------------------------------------------------------

% Margin Definition Badge: Sleek, compact card for margin glossary terms
\NewDocumentCommand{\defmargin}{ m +m }{%
  \marginnote{%
    \begin{tcolorbox}[
      enhanced,
      breakable=false,
      colback=structural!4!white,
      colframe=structural,
      boxrule=0pt,
      leftrule=2pt,
      arc=1pt,
      left=3.5pt,
      right=3pt,
      top=2pt,
      bottom=2.5pt,
      before skip=2pt,
      after skip=3pt
    ]%
      {\fontsize{6.5pt}{8pt}\sffamily\bfseries\scshape\color{structural}Definition}\\[1pt]
      {\fontsize{7.2pt}{9.2pt}\sffamily\bfseries\color{ink}#1}\\[1.5pt]
      {\fontsize{6.8pt}{8.8pt}\sffamily\color{softink}#2}%
    \end{tcolorbox}%
  }%
}

% fullwidth spans the text block plus the margin column, for an artifact the
% reader scans across rather than reads down.
%
%   \begin{fullwidth} ... \end{fullwidth}
%
% Use it sparingly, and only for genuinely wide material. A table of sentences
% does not become legible by getting wider; it needs shorter cells.
\newlength{\fullwidthlen}
\setlength{\fullwidthlen}{\dimexpr\textwidth+\marginparsep+\marginparwidth\relax}
\NewEnviron{fullwidth}{%
  \par\noindent
  \makebox[\textwidth][l]{%
    \begin{minipage}{\fullwidthlen}\BODY\end{minipage}%
  }\par
}


% ---------------------------------------------------------------------------
% 10. Margin Alignment Guides (development aid — toggle)
%
% Draws light frames around the text block, header/footer, and the margin-note
% column, so text, figures, tables, and notes can be aligned during authoring.
% Set \MarginDebugfalse for production builds; \MarginDebugtrue to show guides.
% ---------------------------------------------------------------------------

\newif\ifMarginDebug
\MarginDebugfalse% <<< set to \MarginDebugtrue to show margin alignment guides
\ifMarginDebug
  \usepackage{showframe}
  \renewcommand*\ShowFrameLinethickness{0.2pt}
  \renewcommand*\ShowFrameColor{\color{red!45}}
\fi




% ===========================================================================
% Physical AI Systems, cover and title pages
% ===========================================================================
%
% Loaded after tex/preamble.tex, which defines the palette this file uses.
% This file draws; the preamble configures. Keep it that way.
%
% \maketitle produces two pages:
%
%   1. The cover. Full-bleed ETH blue, crimson rules top and bottom, and the
%      proposal-permission loop drawn as the book's signature diagram. The
%      loop is on the cover because it is the argument: an untrusted proposer,
%      a trusted permitter, an irreversible world, and the feedback that makes
%      the whole thing endogenous.
%   2. The title page. Quiet, typographic, no artwork.
%
% Text on the cover is set here rather than pulled from \@title so the cover
% can be composed as a piece of design. If the book is retitled, change it in
% both places.
%
% ===========================================================================

\makeatletter
\renewcommand*{\maketitle}{%

% --- Page 1: cover ---------------------------------------------------------
\begin{titlepage}\thispagestyle{empty}%
  \newgeometry{margin=0pt}%
  \begin{tikzpicture}[remember picture,overlay]

    % Pure white background
    \fill[white] (current page.south west) rectangle (current page.north east);

    % Floating 3D sculpture artwork centered in upper 58%
    \node[anchor=center, inner sep=0pt] at ([yshift=1.75in]current page.center) {%
      
\includegraphics[width=0.86\paperwidth]{/Users/VJ/GitHub/MLSysBook-integrate-physical/physical/book/assets/covers/cover-physical-loop-recommended-art.png}%
    };

    % Left-aligned Subtitle in Roman Serif (Matching Vol 1 & Vol 2)
    \node[anchor=south west, inner sep=0pt] at ([xshift=0.875in, yshift=3.35in]current page.south west) {%
      {\fontsize{26pt}{30pt}\rmfamily\selectfont\color{softink}Physical AI}%
    };

    % Large Two-Line Title in Book's Native Serif (Matching Vol 1 & Vol 2)
    \node[anchor=north west, inner sep=0pt] at ([xshift=0.875in, yshift=3.15in]current page.south west) {%
      \linespread{1.05}\selectfont%
      {\fontsize{50pt}{56pt}\rmfamily\selectfont\color{ink}%
       \begin{tabular}{@{}l@{}}%
         \fontsize{50pt}{56pt}\rmfamily\selectfont Machine Learning\\[0.10em]%
         \fontsize{50pt}{56pt}\rmfamily\selectfont Systems%
       \end{tabular}%
      }%
    };

    % Right-aligned Author Block at bottom right (Matching Vol 1 & Vol 2)
    \node[anchor=south east, inner sep=0pt] at ([xshift=-0.875in, yshift=0.95in]current page.south east) {%
      \linespread{1.1}\selectfont%
      {\fontsize{24pt}{30pt}\rmfamily\selectfont\color{ink}%
       \begin{tabular}{@{}r@{}}%
         \fontsize{24pt}{30pt}\rmfamily\selectfont Vijay\\[0.15em]%
         \fontsize{24pt}{30pt}\rmfamily\selectfont Janapa Reddi%
       \end{tabular}%
      }%
    };

  \end{tikzpicture}%
  \restoregeometry
\end{titlepage}%

% --- Page 2: interior title page (Matching Volume 1 & Volume 2) ---------
\begin{titlepage}\thispagestyle{empty}%
  \raggedright
  \vspace*{-0.6\topskip}

  {\fontsize{15pt}{17pt}\rmfamily\selectfont Physical AI: Machine Learning Systems\par}
  \vspace{0.8\baselineskip}
  {\fontsize{12pt}{14pt}\rmfamily\itshape\color{accentcolor}That Sense and Act\par}

  \vspace{10\baselineskip}

  {\fontsize{14pt}{16pt}\rmfamily\selectfont Vijay Janapa Reddi\par}
  \vspace{0.4\baselineskip}
  {\fontsize{10pt}{12pt}\rmfamily\selectfont\color{softink}Harvard University \ $\cdot$\  ETH Zurich\par}

  \vfill

  {\fontsize{11pt}{13pt}\rmfamily\selectfont%
   The MIT Press\\[0.4ex]%
   Cambridge, Massachusetts\\[0.4ex]%
   London, England%
  \par}

  \vspace*{0.02\textheight}
\end{titlepage}%

}
\makeatother


\begin{document}
\frontmatter
\maketitle

\tableofcontents

\bookmarksetup{startatroot}

\chapter*{Preface}\label{preface}
\addcontentsline{toc}{chapter}{Preface}

\markboth{Preface}{Preface}

A physical AI system is a machine with three parts.

It has a \textbf{body}, which moves mass through space, applies force
through contact, or regulates a flow of energy. The body obeys mechanics
and thermodynamics rather than instructions. It has momentum that must
be arrested somewhere, windings that accumulate heat faster than they
shed it, and transmissions that absorb whatever discontinuity a command
hands them.

It has a \textbf{brain}, which perceives, remembers, reasons, and plans.
Today that brain is built from learned models, and their value is
precisely that nobody wrote down what they do. A learned component
recognizes an object nobody enumerated and follows an instruction
phrased in a way nobody anticipated. It also arrives at its answer after
a delay that varies by an order of magnitude between one invocation and
the next, and it cannot reliably say when it is wrong.

And it has a \textbf{nervous system}, which carries signals between the
two. It runs on deadlines the body sets rather than the brain, it
decides what reaches an actuator and what does not, and it is the only
part of the machine that can refuse.

The engineering problem this book takes up is that these three parts
have incompatible properties, and they have to work as one machine
anyway. The brain is capable and unaccountable. The body is unforgiving
and cannot be argued with. Between them sits a nervous system that must
be fast enough to matter, predictable enough to trust, and independent
enough to keep working when the brain does not.

\section*{What This Book Is}\label{what-this-book-is}
\addcontentsline{toc}{section}{What This Book Is}

\markright{What This Book Is}

\lettrine{T}{his} is a systems engineering book. It draws on machine
learning, on real-time and embedded systems, and on control, but it is
not a book about any of the three and it does not attempt to teach them.
Its subject is composition: what each part of the machine must guarantee
to the others, how authority over physical action is allocated, what
happens to those guarantees when everything shares one piece of silicon,
and what evidence would justify letting the result operate.

The question underneath every chapter is what the surrounding system
must know, measure, enforce, and be able to show before an unverified
proposal is permitted to move matter.

That question is new. Machine learning systems engineering has spent a
decade on training, serving, and deployment, and it has become a mature
discipline, but almost all of it stops at a digital output. This book
takes up what happens when the output is a current through a gate
driver, and the error is not a wrong label but a sheared gear.

\subsection*{The Three Paradigms}\label{the-three-paradigms}
\addcontentsline{toc}{subsection}{The Three Paradigms}

{\def\LTcaptype{none} % do not increment counter
\begin{longtable}[]{@{}
  >{\raggedright\arraybackslash}p{(\linewidth - 6\tabcolsep) * \real{0.2500}}
  >{\raggedright\arraybackslash}p{(\linewidth - 6\tabcolsep) * \real{0.2500}}
  >{\raggedright\arraybackslash}p{(\linewidth - 6\tabcolsep) * \real{0.2500}}
  >{\raggedright\arraybackslash}p{(\linewidth - 6\tabcolsep) * \real{0.2500}}@{}}
\toprule\noalign{}
\begin{minipage}[b]{\linewidth}\raggedright
Property
\end{minipage} & \begin{minipage}[b]{\linewidth}\raggedright
Traditional Digital ML
\end{minipage} & \begin{minipage}[b]{\linewidth}\raggedright
Classical Robotics \& Control
\end{minipage} & \begin{minipage}[b]{\linewidth}\raggedright
Physical AI
\end{minipage} \\
\midrule\noalign{}
\endhead
\bottomrule\noalign{}
\endlastfoot
\textbf{Observation Stream} & Exogenous, static
(\(x \sim P_{\text{data}}\)). Model errors never alter future test data.
& Dynamic feedback (\(y = Cx\)). Fixed sensor geometry with known noise
models. & \textbf{Endogenous \& Non-Stationary}
(\(s_{t+1} \sim P(\cdot \mid s_t, a_t)\)). Actions shift camera
viewpoints and contact states. \\
\textbf{Error Cost \& Boundary} & Low/Advisory: incorrect tokens or
labels cost compute time or screen pixels. & Bounded by analytical
transfer functions, gain margins, and Lyapunov proofs. &
\textbf{Thermodynamically Irreversible:} Bad actions cost kinetic
collisions, stator thermal damage, and mechanical yield. \\
\textbf{Policy Authority} & High-capacity statistical approximation;
unverified out-of-distribution. & Explicit mathematical formulas (PID,
LQR, H-infinity, analytic MPC). & \textbf{High-capacity learned
proposer} filtered by \textbf{deterministic hardware enforcers}. \\
\textbf{Temporal Deadlines} & Elastic / best-effort
(throughput-optimized batch processing). & Hard real-time cycle
determinism (\(100\text{ Hz}\text{--}1000\text{ Hz}\)). &
\textbf{Multi-Rate Tiered Hierarchy} (\(1000\text{ Hz}\) inner enforcer,
\(10\text{--}100\text{ Hz}\) perception, \(1\text{--}10\text{ Hz}\)
neural policy). \\
\textbf{Safety Certification} & Validation loss, accuracy, perplexity on
i.i.d. splits. & Formal stability proofs and gain/phase margins. &
\textbf{Control Barrier Functions (\(h(x) \ge 0\))}, hardware isolation,
and Claim-Argument-Evidence trees. \\
\end{longtable}
}

\section*{What This Book Is Not}\label{what-this-book-is-not}
\addcontentsline{toc}{section}{What This Book Is Not}

\markright{What This Book Is Not}

It is not a robotics text. Kinematics, rigid-body dynamics, and
manipulator geometry appear only where a systems decision rests on them,
and readers who need those subjects properly should go to Craig, to
Lynch and Park, or to Siciliano.

It is not a controls text. Stability analysis and frequency-domain
methods appear only where they decide an engineering question, and
Slotine and Åström remain the places to learn them.

It is not a machine learning text. No model is discussed apart from what
its use costs in time, memory, energy, and physical consequence, and
Goodfellow, Bishop, and Murphy cover the rest.

It is also not a book about physical AI as the phrase is often used,
which tends to mean generative models applied to robots and treats
physical execution as a downstream detail. This book is about the
engineering that execution demands.

\section*{Who This Book Is For}\label{who-this-book-is-for}
\addcontentsline{toc}{section}{Who This Book Is For}

\markright{Who This Book Is For}

Physical AI is the convergence of three distinct engineering traditions,
and this book is designed so that no reader is left behind and no reader
is talked down to:

\begin{itemize}
\tightlist
\item
  \textbf{The Machine Learning \& AI Practitioner:} You arrive with deep
  mastery of gradient descent, transformer architectures, diffusion
  decoders, and cloud training pipelines. This book builds the physical
  and systems bridge: showing why open-loop statistical accuracy fails
  in dynamic environments, how memory bandwidth limits streaming
  inference, and how to interface unprivileged learned policies with
  hard real-time hardware enforcers.
\item
  \textbf{The Robotics \& Control Systems Engineer:} You arrive with
  rigorous training in rigid-body dynamics, differential kinematics,
  Kalman filtering, and Lyapunov stability. This book shows how to
  augment analytical state pipelines with high-capacity learned
  foundation models, replacing fragile handcrafted feature detectors
  with open-vocabulary spatial perception while preserving formal
  Control Barrier Function guarantees.
\item
  \textbf{The Embedded Systems \& Silicon Architect:} You arrive with
  command over bare-metal C, real-time microcontrollers, interrupt
  service routines, bus protocols, and zero-allocation determinism. This
  book shows how to architect heterogeneous System-on-Chip (SoC)
  platforms where Linux-based neural accelerators and hard real-time
  safety cores share silicon, arbitrating memory crossbars without
  starving physical safety reflexes.
\end{itemize}

Every unfamiliar mechanism in this book is introduced first by the
\textbf{physical phenomenon it governs}, second by its \textbf{systems
interface}, and only third by its \textbf{disciplinary name}---ensuring
that engineers from any background build shared physical and
architectural intuition from first principles.

The book is written for advanced undergraduates, graduate students, and
practicing systems engineers, assuming basic familiarity with calculus,
linear algebra, probability, and standard programming, but requiring no
prior robotics or embedded systems background.

\section*{How the Book Is Organized}\label{how-the-book-is-organized}
\addcontentsline{toc}{section}{How the Book Is Organized}

\markright{How the Book Is Organized}

The book is organized into four sequential parts:

\begin{itemize}
\tightlist
\item
  \textbf{Part I: Physical AI Foundations (Chapters 1--4):} Establishes
  the physical container. Chapter 1 introduces the causal boundary and
  the three canonical machine archetypes. Chapter 2 derives the
  non-negotiable physical limits of the mechanical body. Chapter 3
  formalizes the neural brain as a high-capacity, unprivileged compute
  component under memory bandwidth walls. Chapter 4 derives the nervous
  system's multi-rate execution clocks, lock-free memory transport, and
  the master 5x5 synthesis grid.
\item
  \textbf{Part II: Data and Learning Systems (Chapters 5--7):} Opens the
  policy lifecycle. Chapter 5 accounts for the collection,
  teleoperation, timestamp synchronization, and compounding error of
  physical demonstration data. Chapter 6 examines the training systems
  trade-offs between imitation learning (ACT, Diffusion Policy),
  reinforcement learning, and sim-to-real physics. Chapter 7 establishes
  the systems limits of offline evaluation and constructs closed-loop
  physical benchmarking protocols.
\item
  \textbf{Part III: The Runtime Embodied Stack (Chapters 8--13):} Walks
  the runtime pipeline stage by stage: Perception Systems (Chapter 8),
  Spatial Memory \& State (Chapter 9), Intent \& Task Reasoning (Chapter
  10), Trajectory Planning (Chapter 11), and Real-Time Reflexes (Chapter
  12). Chapter 13 co-locates stochastic neural accelerators and hard
  real-time safety enforcers on shared silicon, proving when temporal
  and memory isolation holds.
\item
  \textbf{Part IV: Integration, Verification, and Deployment (Chapters
  14--17):} Assembles the systems engineering discipline for shipping
  and operating physical machines. Chapter 14 formalizes shared
  autonomy, human intervention authority, and bumpless transfer. Chapter
  15 ascends the qualification ladder through hardware-in-the-loop (HIL)
  testing and synthetic fault injection. Chapter 16 structures
  production release engineering, telemetry monitoring, and over-the-air
  updates. Chapter 17 explores open frontier challenges across
  foundation models and physical systems.
\end{itemize}

\section*{A Note on What Is Not Yet
Known}\label{a-note-on-what-is-not-yet-known}
\addcontentsline{toc}{section}{A Note on What Is Not Yet Known}

\markright{A Note on What Is Not Yet Known}

A book that teaches engineers to demand evidence cannot itself
overclaim. Where a number in these pages was measured, the measurement
is described. Where it was derived, the derivation is shown. Where it is
illustrative, it says so. And where the honest answer is that the field
does not yet know, the chapter says that instead of reaching for a
citation that does not settle it.

The separation of proposal from permission, which runs through the whole
book, is a teaching architecture and an engineering discipline. It is
not a claim of certified functional safety, and nothing here should be
read as a certification method. What the book teaches is the reasoning
that has to happen before any such process begins.

\bookmarksetup{startatroot}

\chapter*{Course Syllabus: Physical AI Systems}\label{sec-syllabus}
\addcontentsline{toc}{chapter}{Course Syllabus: Physical AI Systems}

\markboth{Course Syllabus: Physical AI Systems}{Course Syllabus:
Physical AI Systems}

\textbf{Project Seminar \& Hardware Studio} · Harvard University · ETH
Zurich

{\def\LTcaptype{none} % do not increment counter
\begin{longtable}[]{@{}
  >{\raggedright\arraybackslash}p{(\linewidth - 2\tabcolsep) * \real{0.5000}}
  >{\raggedright\arraybackslash}p{(\linewidth - 2\tabcolsep) * \real{0.5000}}@{}}
\toprule\noalign{}
\begin{minipage}[b]{\linewidth}\raggedright
Parameter
\end{minipage} & \begin{minipage}[b]{\linewidth}\raggedright
Specification
\end{minipage} \\
\midrule\noalign{}
\endhead
\bottomrule\noalign{}
\endlastfoot
\textbf{Course} & Physical AI Systems (Harvard CS/EE 288 · ETH Zurich
151-0851) \\
\textbf{Credits} & 6 ECTS / 4 Harvard Credits
\emph{(\textasciitilde150--180 total student hours)} \\
\textbf{Format} & Weekly Interactive Seminar + Hands-On Hardware Studio
· \textbf{No Written Exam} \\
\textbf{Language} & English \\
\textbf{Level} & Advanced Undergraduate \& Graduate (MSc / PhD) \\
\textbf{Hardware Kit} & \textbf{Physical AI Kit} (Arduino UNO Q
Dual-Brain: Linux MPU + Real-Time MCU) \\
\textbf{Textbook} & \emph{Physical AI: Machine Learning Systems That
Sense and Act} \\
\textbf{Course Portal} & \url{https://physical.mlsysbook.ai} \\
\end{longtable}
}

\begin{center}\rule{0.5\linewidth}{0.5pt}\end{center}

\section*{Pedagogical Vision}\label{pedagogical-vision}
\addcontentsline{toc}{section}{Pedagogical Vision}

\markright{Pedagogical Vision}

\lettrine{C}{lassical} machine learning curricula operate entirely
\textbf{behind glass}: a misplaced class label or high prediction error
results in a software retry or log entry. \textbf{Physical AI begins
when software commands physical actuators}---mass moves, friction
generates heat, kinetic energy is spent, and the physical world
permanently transitions (\(s_{t+1} \sim P(s \mid s_t, u_t)\)). You
cannot \texttt{ctrl+z} kinetic energy.

This course teaches the systems engineering disciplines required for
that transition:
\[\text{Observe} \longrightarrow \text{Estimate} \longrightarrow \text{Intend} \longrightarrow \text{Plan} \longrightarrow \text{Enforce}\]

The central invariant governing the course is the
\textbf{Proposal--Permission Boundary}: 1. \textbf{The Cognitive Cortex
(Linux MPU):} Runs deep multimodal foundation models (VLMs), spatial
tokenizers (DINOv2), and generative action chunking decoders (Diffusion
Policy / ACT). It generates timestamped, expiring intent proposals. 2.
\textbf{The Safety Enforcer (Real-Time MCU):} Runs a bare-metal,
zero-allocation \(1000\text{ Hz}\) control loop. It verifies Control
Barrier Functions (\(h(x) \ge 0\)), monitors dynamic stopping distances
(\(d_{\text{stop}}\)), and commands motor gate inverters.

\subsection*{Multi-Disciplinary Curriculum
Pathways}\label{multi-disciplinary-curriculum-pathways}
\addcontentsline{toc}{subsection}{Multi-Disciplinary Curriculum
Pathways}

Because students enter this curriculum from diverse engineering
backgrounds, the course provides three tailored conceptual pathways:

\begin{itemize}
\tightlist
\item
  \textbf{Pathway A · The Machine Learning Practitioner:} Emphasizes the
  thermodynamic and kinematic consequences of model decisions. Focuses
  on why open-loop statistical accuracy fails in dynamic environments,
  how memory bandwidth bounds inference frequency
  (\(t_{\text{stream}} = M_{\text{weights}} / B_{\text{mem}}\)), why
  compounding covariate drift requires closed-loop exposure bounds, and
  how to interface unprivileged neural models with deterministic
  real-time supervisors (Key Chapters: Chapter~\ref{sec-boundary},
  Chapter~\ref{sec-body}, Chapter~\ref{sec-brain},
  Chapter~\ref{sec-data}, Chapter~\ref{sec-training},
  Chapter~\ref{sec-intent}, Chapter~\ref{sec-enforcement},
  Chapter~\ref{sec-release}).
\item
  \textbf{Pathway B · The Robotics \& Control Engineer:} Emphasizes
  augmenting classical analytical dynamics with high-capacity learned
  foundation models. Focuses on open-vocabulary 3D semantic grounding,
  learning multimodal action distributions via trajectory diffusion
  without mode collapse, and enforcing Control Barrier Function (CBF-QP)
  safety shields over learned trajectory proposals (Key Chapters:
  Chapter~\ref{sec-boundary}, Chapter~\ref{sec-brain},
  Chapter~\ref{sec-nervous}, Chapter~\ref{sec-perception},
  Chapter~\ref{sec-memory}, Chapter~\ref{sec-planning},
  Chapter~\ref{sec-enforcement}, Chapter~\ref{sec-verification}).
\item
  \textbf{Pathway C · The Embedded \& Silicon Architect:} Emphasizes
  heterogeneous system partitioning, hardware arbitration, and
  deterministic safety enforcement. Focuses on multi-rate clock
  cadences, lock-free seqlocks with memory barriers, AXI QoS bus
  arbitration under neural DMA bursts, inductive voltage droop
  (\(L \, dI/dt\)), and hardware fault injection (Key Chapters:
  Chapter~\ref{sec-boundary}, Chapter~\ref{sec-body},
  Chapter~\ref{sec-nervous}, Chapter~\ref{sec-memory},
  Chapter~\ref{sec-enforcement}, Chapter~\ref{sec-placement},
  Chapter~\ref{sec-intervention}, Chapter~\ref{sec-verification}).
\end{itemize}

\begin{center}\rule{0.5\linewidth}{0.5pt}\end{center}

\section*{The 14-Week Curriculum \& Laboratory
Spine}\label{the-14-week-curriculum-laboratory-spine}
\addcontentsline{toc}{section}{The 14-Week Curriculum \& Laboratory
Spine}

\markright{The 14-Week Curriculum \& Laboratory Spine}

The course is structured across 17 textbook chapters and 14 laboratory
studio sessions:

\subsection*{Part I: The Laws of Physical Action (Weeks
1--4)}\label{part-i-the-laws-of-physical-action-weeks-14}
\addcontentsline{toc}{subsection}{Part I: The Laws of Physical Action
(Weeks 1--4)}

\begin{itemize}
\tightlist
\item
  \textbf{Week 1 · Kit Bring-Up \& The Closed Loop
  (Chapter~\ref{sec-boundary}):}

  \begin{itemize}
  \tightlist
  \item
    \emph{Lecture:} The causal loop, state mutation, and the
    first-exhausted physical budget.
  \item
    \emph{Studio Lab (\texttt{00-kit-bringup},
    \texttt{01-close-the-loop}):} Dual-brain IPC setup, safe idle
    transitions, open-loop vs closed-loop telemetry.
  \item
    \emph{Deliverable:} \textbf{Loop Charter}.
  \end{itemize}
\item
  \textbf{Week 2 · Reflected Inertia \& Actuator Dynamics
  (Chapter~\ref{sec-body}):}

  \begin{itemize}
  \tightlist
  \item
    \emph{Lecture:} Reflected rotor inertia (\(N^2 J_{\text{rotor}}\)),
    QDD vs.~Harmonic drives, and thermal impedance networks
    (\(\Delta T = P \cdot \theta_{\text{JA}}\)).
  \item
    \emph{Studio Lab (\texttt{02-freshness-wall}):} Inverter current
    shunt measurement, thermal step response, stopping distance bounds.
  \end{itemize}
\item
  \textbf{Week 3 · The Two-Speed Brain \& Multi-Rate Pipelines
  (Chapter~\ref{sec-brain}, Chapter~\ref{sec-nervous}):}

  \begin{itemize}
  \tightlist
  \item
    \emph{Lecture:} Slow deliberative foundation models
    (\(1\text{--}5\text{ Hz}\)) vs.~fast reflex enforcers
    (\(1000\text{ Hz}\)); lock-free seqlocks and zero-allocation SRAM.
  \item
    \emph{Studio Lab (\texttt{03-measure-both-brains},
    \texttt{04-runtime-fault-containment}):} \(P_{99}\) latency
    profiling, inter-processor seqlock buffers, MPU crash survival.
  \item
    \emph{Deliverable:} \textbf{Requirements Ledger}.
  \end{itemize}
\item
  \textbf{Week 4 · Data Foundations \& Teleoperation
  (Chapter~\ref{sec-data}):}

  \begin{itemize}
  \tightlist
  \item
    \emph{Lecture:} Compounding covariate shift, distribution drift,
    bilateral teleoperation latency, and Open X-Embodiment datasets.
  \item
    \emph{Studio Lab (\texttt{04-dma-tokens}):} Human demonstration
    capture, teleoperation jitter compensation, action space
    normalization.
  \end{itemize}
\end{itemize}

\subsection*{Part II: The Architectural Spine (Weeks
5--8)}\label{part-ii-the-architectural-spine-weeks-58}
\addcontentsline{toc}{subsection}{Part II: The Architectural Spine
(Weeks 5--8)}

\begin{itemize}
\tightlist
\item
  \textbf{Week 5 · Perception \& 3D Spatial Encoders
  (Chapter~\ref{sec-perception}):}

  \begin{itemize}
  \tightlist
  \item
    \emph{Lecture:} Sensor latency waterfalls, rolling shutter
    distortion, on-manifold IMU preintegration, and open-vocabulary 3D
    scene graphs.
  \item
    \emph{Studio Lab (\texttt{05-perception-frontier}):} MIPI CSI-2 DMA
    ring buffers, camera-to-IMU temporal calibration, spatial feature
    tokenization.
  \item
    \emph{Deliverable:} \textbf{Observation Contract}.
  \end{itemize}
\item
  \textbf{Week 6 · Spatial Memory \& Epistemic Drift
  (Chapter~\ref{sec-memory}):}

  \begin{itemize}
  \tightlist
  \item
    \emph{Lecture:} \(SE(3)\) coordinate frame trees, volumetric
    raycasting (3D-DDA), TSDFs, and covariance expansion.
  \item
    \emph{Studio Lab (\texttt{06-belief-drift}):} Timestamp skew
    mitigation, occupancy voxel updates, belief lease expiration.
  \item
    \emph{Deliverable:} \textbf{State \& Timing Model}.
  \end{itemize}
\item
  \textbf{Week 7 · Intent \& Open-Vocabulary Grounding
  (Chapter~\ref{sec-intent}):}

  \begin{itemize}
  \tightlist
  \item
    \emph{Lecture:} Vision-Language Models (VLMs) as unprivileged
    proposal engines, spatial affordances, and expiring intent leases.
  \item
    \emph{Studio Lab (\texttt{07-two-speed-intent}):} Edge VLM prompt
    parsing, 3D target bounding, intent lease TTL enforcement.
  \item
    \emph{Deliverable:} \textbf{Intent Schema}.
  \end{itemize}
\item
  \textbf{Week 8 · The Signature Enforcer \& Midterm Milestone
  (Chapter~\ref{sec-enforcement}):}

  \begin{itemize}
  \tightlist
  \item
    \emph{Lecture:} Control Barrier Functions (CBFs), active-set barrier
    quadratic programming (QP), and minimal intervention safety shields.
  \item
    \emph{Studio Lab (\texttt{08-mcu-enforcer}):} \textbf{Signature
    Lab:} Bare-metal \(1000\text{ Hz}\) MCU CBF veto filter overriding
    invalid MPU trajectory proposals.
  \item
    \emph{Deliverable:} \textbf{Enforcement Design}.
  \end{itemize}
\end{itemize}

\subsection*{Part III: Realization, Governance \& Release (Weeks
9--14)}\label{part-iii-realization-governance-release-weeks-914}
\addcontentsline{toc}{subsection}{Part III: Realization, Governance \&
Release (Weeks 9--14)}

\begin{itemize}
\tightlist
\item
  \textbf{Week 9 · Heterogeneous Silicon Placement
  (Chapter~\ref{sec-placement}):}

  \begin{itemize}
  \tightlist
  \item
    \emph{Lecture:} AXI crossbar contention, memory bandwidth walls,
    DRAM bank conflicts, and PDN inductive voltage droop
    (\(\Delta V = L \cdot dI/dt\)).
  \item
    \emph{Studio Lab (\texttt{09-placement-ripple}):} Memory bandwidth
    throttling, bus stress injection, real-time core isolation.
  \item
    \emph{Deliverable:} \textbf{Placement Ledger}.
  \end{itemize}
\item
  \textbf{Week 10 · Human Intervention \& Shared Autonomy
  (Chapter~\ref{sec-intervention}):}

  \begin{itemize}
  \tightlist
  \item
    \emph{Lecture:} Shared control policy blending, bumpless control
    transfer, and human takeover reaction times.
  \item
    \emph{Studio Lab (\texttt{10-shadow-and-faults},
    \texttt{11-authority-paths}):} Bumpless joystick override,
    state-conditioned authority arbitration.
  \item
    \emph{Deliverable:} \textbf{Authority Design}.
  \end{itemize}
\item
  \textbf{Week 11 · Verification \& Synthetic Falsification
  (Chapter~\ref{sec-verification}):}

  \begin{itemize}
  \tightlist
  \item
    \emph{Lecture:} The 4-stage qualification ladder (Sim \(\to\) PIL
    \(\to\) HIL \(\to\) In-Situ), temporal logic falsification, and
    hardware-in-the-loop stress testing.
  \item
    \emph{Studio Lab (\texttt{12-learning-turn},
    \texttt{13-ship-gate}):} Cross-layer fault injection (dropped
    packets, voltage sags, optical occlusion).
  \end{itemize}
\item
  \textbf{Week 12 · Release Assurance \& Safety Cases
  (Chapter~\ref{sec-release}):}

  \begin{itemize}
  \tightlist
  \item
    \emph{Lecture:} Claim-Argument-Evidence (CAE) safety cases, Goal
    Structuring Notation (GSN), and UL 4600 / ISO standard alignment.
  \item
    \emph{Studio Lab (\texttt{99-design-review} prep):} Cryptographic
    artifact signing, safety case construction, release gate audit.
  \item
    \emph{Deliverable:} \textbf{Release Case}.
  \end{itemize}
\item
  \textbf{Weeks 13--14 · Capstone Jury \& Unannounced Fault Defense
  (Chapter~\ref{sec-frontier}):}

  \begin{itemize}
  \tightlist
  \item
    \emph{Studio Defense:} Live capstone demonstration. The teaching
    team injects unannounced real-time hardware/sensor faults. Teams
    must prove that their physical agent either safely achieves the task
    or transitions into a certified safe holding state without
    mechanical impact.
  \item
    \emph{Deliverable:} \textbf{Final Signed Release Verdict}.
  \end{itemize}
\end{itemize}

\begin{center}\rule{0.5\linewidth}{0.5pt}\end{center}

\section*{Grading \& Assessment Policy}\label{grading-assessment-policy}
\addcontentsline{toc}{section}{Grading \& Assessment Policy}

\markright{Grading \& Assessment Policy}

This course employs a \textbf{Project-First, Continuous-Evidence
Evaluation} model without classical written exams:

{\def\LTcaptype{none} % do not increment counter
\begin{longtable}[]{@{}
  >{\raggedright\arraybackslash}p{(\linewidth - 4\tabcolsep) * \real{0.3077}}
  >{\centering\arraybackslash}p{(\linewidth - 4\tabcolsep) * \real{0.3846}}
  >{\raggedright\arraybackslash}p{(\linewidth - 4\tabcolsep) * \real{0.3077}}@{}}
\toprule\noalign{}
\begin{minipage}[b]{\linewidth}\raggedright
Component
\end{minipage} & \begin{minipage}[b]{\linewidth}\centering
Weight
\end{minipage} & \begin{minipage}[b]{\linewidth}\raggedright
Assessment Criteria
\end{minipage} \\
\midrule\noalign{}
\endhead
\bottomrule\noalign{}
\endlastfoot
\textbf{Studio Labs \& Checkpoints} & 25\% & Timely execution of weekly
bench milestones, hardware wiring integrity, and working firmware
builds. \\
\textbf{Engineering Notebook} & 35\% & Rigorous, version-controlled
records of latency distributions (\(P_{50}, P_{90}, P_{99}\)), energy
budgets, and system logs across all chapter milestones. \\
\textbf{Midterm Enforcement Gate} & 15\% & Formal verification and live
bench demonstration of the Lab 08 1 kHz MCU Control Barrier Function
enforcer. \\
\textbf{Capstone System Defense} & 25\% & Live final demonstration
defending the Physical Agent against unannounced seeded faults before
the faculty jury. \\
\end{longtable}
}

\bookmarksetup{startatroot}

\chapter*{Mathematical and Systems Notation}\label{sec-notation}
\addcontentsline{toc}{chapter}{Mathematical and Systems Notation}

\markboth{Mathematical and Systems Notation}{Mathematical and Systems
Notation}

To ensure complete clarity and consistency across all seventeen
chapters, this textbook adopts a unified mathematical, computational,
and physical systems notation.

\begin{center}\rule{0.5\linewidth}{0.5pt}\end{center}

\section*{1. Physical State and
Dynamics}\label{physical-state-and-dynamics}
\addcontentsline{toc}{section}{1. Physical State and Dynamics}

\markright{1. Physical State and Dynamics}

{\def\LTcaptype{none} % do not increment counter
\begin{longtable}[]{@{}
  >{\raggedright\arraybackslash}p{(\linewidth - 4\tabcolsep) * \real{0.3333}}
  >{\raggedright\arraybackslash}p{(\linewidth - 4\tabcolsep) * \real{0.3333}}
  >{\raggedright\arraybackslash}p{(\linewidth - 4\tabcolsep) * \real{0.3333}}@{}}
\toprule\noalign{}
\begin{minipage}[b]{\linewidth}\raggedright
Symbol
\end{minipage} & \begin{minipage}[b]{\linewidth}\raggedright
Definition
\end{minipage} & \begin{minipage}[b]{\linewidth}\raggedright
Canonical Units
\end{minipage} \\
\midrule\noalign{}
\endhead
\bottomrule\noalign{}
\endlastfoot
\(s_t \in \mathcal{S}\) & True physical state of the machine and
environment at time \(t\) & Mixed / Continuous \\
\(o_t \in \mathcal{O}\) & Sensor observation ingested by the computing
pipeline & Pixels / Volts / Rad \\
\(a_t \in \mathcal{A}\) & Actuator action command (e.g., target torque
or joint setpoint) & \(\text{N}\cdot\text{m}\) / \(\text{rad}\) \\
\(\mathbf{q} \in \mathbb{R}^n\) & Generalized joint coordinate
configuration vector & \(\text{rad}\) or \(\text{m}\) \\
\(\dot{\mathbf{q}} \in \mathbb{R}^n\) & Generalized joint velocity
vector & \(\text{rad/s}\) or \(\text{m/s}\) \\
\(\ddot{\mathbf{q}} \in \mathbb{R}^n\) & Generalized joint acceleration
vector & \(\text{rad/s}^2\) or \(\text{m/s}^2\) \\
\(\boldsymbol{\tau} \in \mathbb{R}^n\) & Generalized joint actuator
torque/force vector & \(\text{N}\cdot\text{m}\) or \(\text{N}\) \\
\(\mathbf{M}(\mathbf{q})\) & Generalized inertia / mass matrix &
\(\text{kg}\cdot\text{m}^2\) \\
\(\mathbf{C}(\mathbf{q}, \dot{\mathbf{q}})\) & Coriolis and centrifugal
force matrix & \(\text{N}\cdot\text{m}\cdot\text{s/rad}\) \\
\(\mathbf{g}(\mathbf{q})\) & Gravitational generalized force vector &
\(\text{N}\cdot\text{m}\) \\
\(\mathbf{J}(\mathbf{q})\) & Kinematic manipulator Jacobian matrix &
\(\text{m/rad}\) \\
\(N\) & Transmission / gearbox reduction ratio (\(N \gg 1\)) &
Dimensionless \\
\(N^2 J\) & Reflected rotor inertia seen at the gearbox output shaft &
\(\text{kg}\cdot\text{m}^2\) \\
\end{longtable}
}

\begin{center}\rule{0.5\linewidth}{0.5pt}\end{center}

\section*{2. Geometry and Spatial
Transformations}\label{geometry-and-spatial-transformations}
\addcontentsline{toc}{section}{2. Geometry and Spatial Transformations}

\markright{2. Geometry and Spatial Transformations}

{\def\LTcaptype{none} % do not increment counter
\begin{longtable}[]{@{}
  >{\raggedright\arraybackslash}p{(\linewidth - 4\tabcolsep) * \real{0.3333}}
  >{\raggedright\arraybackslash}p{(\linewidth - 4\tabcolsep) * \real{0.3333}}
  >{\raggedright\arraybackslash}p{(\linewidth - 4\tabcolsep) * \real{0.3333}}@{}}
\toprule\noalign{}
\begin{minipage}[b]{\linewidth}\raggedright
Symbol
\end{minipage} & \begin{minipage}[b]{\linewidth}\raggedright
Definition
\end{minipage} & \begin{minipage}[b]{\linewidth}\raggedright
Canonical Dimension
\end{minipage} \\
\midrule\noalign{}
\endhead
\bottomrule\noalign{}
\endlastfoot
\(SE(3)\) & Special Euclidean Group of 3D rigid-body poses
(\(\mathbb{R}^3 \times SO(3)\)) & Lie Group \\
\(SO(3)\) & Special Orthogonal Group of 3D rotational matrices &
\(3 \times 3\) Orthogonal \\
\(\mathbf{p} \in \mathbb{R}^3\) & Translation position vector
\([x, y, z]^\top\) & \(\text{m}\) \\
\(\mathbf{R} \in SO(3)\) & 3D rotation matrix
(\(\mathbf{R}^\top \mathbf{R} = \mathbf{I}, \det(\mathbf{R}) = +1\)) &
\(3 \times 3\) Matrix \\
\(\mathbf{q}_{\text{quat}} \in \mathbb{H}\) & Unit quaternion rotation
representation \([w, x, y, z]^\top\) &
\(\|\mathbf{q}_{\text{quat}}\| = 1\) \\
\(\mathbf{T}_{A \to B}\) & Homogeneous rigid transform mapping
coordinates from Frame \(A\) to Frame \(B\) & \(4 \times 4\) Matrix \\
\(\boldsymbol{\Sigma}_{\text{depth}}\) & 3D spatial covariance matrix of
perception depth estimate & \(\text{m}^2\) \\
\end{longtable}
}

\begin{center}\rule{0.5\linewidth}{0.5pt}\end{center}

\section*{3. Cognitive Deliberation and Policy
Models}\label{cognitive-deliberation-and-policy-models}
\addcontentsline{toc}{section}{3. Cognitive Deliberation and Policy
Models}

\markright{3. Cognitive Deliberation and Policy Models}

{\def\LTcaptype{none} % do not increment counter
\begin{longtable}[]{@{}
  >{\raggedright\arraybackslash}p{(\linewidth - 4\tabcolsep) * \real{0.3333}}
  >{\raggedright\arraybackslash}p{(\linewidth - 4\tabcolsep) * \real{0.3333}}
  >{\raggedright\arraybackslash}p{(\linewidth - 4\tabcolsep) * \real{0.3333}}@{}}
\toprule\noalign{}
\begin{minipage}[b]{\linewidth}\raggedright
Symbol
\end{minipage} & \begin{minipage}[b]{\linewidth}\raggedright
Definition
\end{minipage} & \begin{minipage}[b]{\linewidth}\raggedright
Canonical Units
\end{minipage} \\
\midrule\noalign{}
\endhead
\bottomrule\noalign{}
\endlastfoot
\(\pi_\theta(a \mid o)\) & Parametric neural policy with weights
\(\theta\) & Probability density \\
\(\mathbf{z} \in \mathbb{R}^d\) & Latent spatial-semantic representation
vector & Dimensionless \\
\(H\) & Action chunk temporal prediction horizon (e.g., \(H=16\) or
\(H=32\)) & Discrete time steps \\
\(\mathbf{A}_{t:t+H}\) & Predicted action chunk trajectory vector
\([a_t, a_{t+1}, \dots, a_{t+H-1}]\) & Sequence of actions \\
\(\mathcal{D}_{\text{demo}}\) & Dataset of nominal human/expert
demonstrations & Sequence of tuples \\
\(\mathcal{D}_{\text{rec}}\) & Dataset of human recovery and
intervention trajectories & Sequence of tuples \\
\(\epsilon_{\text{BC}}\) & Single-step imitation learning policy error
rate & Dimensionless (\(\le 1\)) \\
\(O(T^2 \epsilon)\) & Compounding quadratic covariate drift error under
open-loop policy rollouts & Cumulative error \\
\(\mathcal{L}_{\text{intent}}\) & Expiring Intent Lease with bounded
Time-To-Live (TTL) & Spatial-temporal contract \\
\end{longtable}
}

\begin{center}\rule{0.5\linewidth}{0.5pt}\end{center}

\section*{4. Real-Time Control, Safety, and
Enforcement}\label{real-time-control-safety-and-enforcement}
\addcontentsline{toc}{section}{4. Real-Time Control, Safety, and
Enforcement}

\markright{4. Real-Time Control, Safety, and Enforcement}

{\def\LTcaptype{none} % do not increment counter
\begin{longtable}[]{@{}
  >{\raggedright\arraybackslash}p{(\linewidth - 4\tabcolsep) * \real{0.3333}}
  >{\raggedright\arraybackslash}p{(\linewidth - 4\tabcolsep) * \real{0.3333}}
  >{\raggedright\arraybackslash}p{(\linewidth - 4\tabcolsep) * \real{0.3333}}@{}}
\toprule\noalign{}
\begin{minipage}[b]{\linewidth}\raggedright
Symbol
\end{minipage} & \begin{minipage}[b]{\linewidth}\raggedright
Definition
\end{minipage} & \begin{minipage}[b]{\linewidth}\raggedright
Canonical Units
\end{minipage} \\
\midrule\noalign{}
\endhead
\bottomrule\noalign{}
\endlastfoot
\(\mathcal{C} \subset \mathcal{S}\) & Forward-invariant safe operating
set & State space subset \\
\(h(\mathbf{x}) \ge 0\) & Continuously differentiable Control Barrier
Function (CBF) & Real scalar \\
\(L_f h(\mathbf{x}), L_g h(\mathbf{x})\) & Lie derivatives of the
barrier function along drift \(f\) and control \(g\) & Real scalars \\
\(\gamma\) & Extended class-\(\mathcal{K}\) barrier decay rate parameter
& \(\text{s}^{-1}\) \\
\(u_{\text{nom}}\) & Unprivileged nominal control action proposed by
neural policy & \(\text{N}\cdot\text{m}\) \\
\(u^*\) & Verified, minimally-intervened control action emitted by
CBF-QP filter & \(\text{N}\cdot\text{m}\) \\
\(d_{\text{stop}}\) & Dynamic stopping distance envelope &
\(\text{m}\) \\
\(v_{\max}, a_{\max}, j_{\max}\) & Maximum physical velocity,
acceleration, and jerk limits &
\(\text{m/s}, \text{m/s}^2, \text{m/s}^3\) \\
\end{longtable}
}

\begin{center}\rule{0.5\linewidth}{0.5pt}\end{center}

\section*{5. Hardware Silicon and Electrical
Budgets}\label{hardware-silicon-and-electrical-budgets}
\addcontentsline{toc}{section}{5. Hardware Silicon and Electrical
Budgets}

\markright{5. Hardware Silicon and Electrical Budgets}

{\def\LTcaptype{none} % do not increment counter
\begin{longtable}[]{@{}
  >{\raggedright\arraybackslash}p{(\linewidth - 4\tabcolsep) * \real{0.3333}}
  >{\raggedright\arraybackslash}p{(\linewidth - 4\tabcolsep) * \real{0.3333}}
  >{\raggedright\arraybackslash}p{(\linewidth - 4\tabcolsep) * \real{0.3333}}@{}}
\toprule\noalign{}
\begin{minipage}[b]{\linewidth}\raggedright
Symbol
\end{minipage} & \begin{minipage}[b]{\linewidth}\raggedright
Definition
\end{minipage} & \begin{minipage}[b]{\linewidth}\raggedright
Canonical Units
\end{minipage} \\
\midrule\noalign{}
\endhead
\bottomrule\noalign{}
\endlastfoot
\(B_{\text{mem}}\) & Main system memory bus streaming bandwidth &
\(\text{GB/s}\) \\
\(M_{\text{weights}}\) & Neural policy total parameter model memory
footprint & \(\text{GB}\) \\
\(t_{\text{stream}}\) & Minimum memory bandwidth parameter streaming
time (\(M_{\text{weights}}/B_{\text{mem}}\)) & \(\text{s}\) or
\(\text{ms}\) \\
\(\Delta t_{\text{fresh}}\) & Physical freshness deadline for sensor
ingestion & \(\text{ms}\) \\
\(t_{\text{lag}}\) & End-to-end sense-to-torque physical execution
latency & \(\text{ms}\) \\
\(f_{\text{ctrl}}\) & Inner reflex loop execution frequency (canonical:
1000 Hz) & \(\text{Hz}\) \\
\(I\) & Motor phase electrical current & \(\text{Amperes (A)}\) \\
\(R\) & Electrical phase winding or power bus resistance &
\(\text{Ohms } (\Omega)\) \\
\(C_{\text{th}}\) & Stator thermal capacitance & \(\text{J/K}\) \\
\(\Delta T\) & Joulean stator thermal rise
(\(\frac{I^2 R \Delta t}{C_{\text{th}}}\)) & \(\text{Kelvin (K)}\) \\
\(\Delta V\) & Inductive power rail dynamic droop
(\(L \frac{dI}{dt} + IR\)) & \(\text{Volts (V)}\) \\
\end{longtable}
}

\bookmarksetup{startatroot}

\chapter{\texorpdfstring{The Causal
Boundary\chaptersubtitle{Introduction to Physical AI Systems}}{The Causal Boundary}}\label{sec-boundary}

\begin{figure}[H]

\centering{

\includegraphics[width=1\linewidth,height=\textheight,keepaspectratio]{/Users/VJ/GitHub/MLSysBook-integrate-physical/physical/book//Users/VJ/GitHub/MLSysBook-integrate-physical/physical/book/chapters/01-boundary/figures/fig_locator.pdf}

}

\caption{\label{fig-locator}\textbf{Physical AI Systems Architecture
Stack.} The master systems hierarchy spanning System Governance (Layer
4), Cognitive Deliberation (Layer 3), Hard Real-Time Reflex \& Silicon
(Layer 2), and the Physical Body (Layer 1). In this introductory
chapter, our active focus centers on the causal boundary at the
interface between digital software and the physical plant.}

\end{figure}%

Across humanoid robotics, autonomous transport fleets, and surgical
workcells, high-capacity neural foundation models are stepping out of
cloud data centers to command physical machines. Yet inside engineering
design reviews at companies from Tesla and Figure to Boston Dynamics and
Waymo, teams find themselves locked in an architectural crisis: machine
learning researchers celebrate 99\% virtual benchmark scores, while
robotics engineers watch hardware shake, overheat, and strip precision
gearboxes on the physical floor. An error that costs nothing on a screen
destroys hardware in the physical world.

This book is the founding curriculum for \textbf{Physical AI Systems
Architecture}---the discipline of embedding high-capacity learned
statistical models into closed-loop machines governed by classical
mechanics, thermodynamics, and hard real-time silicon determinism. It
bridges open-world cognitive intelligence and dynamical physics,
teaching engineers how to build dependable, certifiable machines that
sense, reason, and act in the physical world.

\newpage{}

\protect\phantomsection\label{callout-objectiveux2a-1.1}
\begin{fbxSimple}{callout-objective}{Learning Objectives}{}
\phantomsection\label{callout-objective*-1.1}

After studying this chapter, you will be able to:

\begin{enumerate}
\def\labelenumi{\arabic{enumi}.}
\tightlist
\item
  Locate the causal boundary in a system architecture and state what
  changes on each side of it.
\item
  Diagnose the ``Tower of Babel'' conflict between ML, Systems,
  Controls, and Embedded disciplines.
\item
  Classify an embodied machine into one of three canonical physical
  archetypes and identify which physical budget binds it first.
\item
  Apply the three-pillar scope test to determine whether a given system
  requires the engineering discipline of physical AI.
\item
  Identify the Four Pillars of an embodied machine: Dynamical Body,
  Cognitive Brain, Real-Time Nervous System, and Governance Envelope.
\item
  Trace how latency, classification instability, and missing authority
  boundaries caused the 2018 Tempe autonomous vehicle collision.
\item
  State the Four Bedrock Laws that govern physical machines across the
  causal boundary.
\item
  Explain how the heterogeneous Dual-Brain architecture (MPU proposer +
  MCU referee across shared SRAM) enforces physical safety invariants.
\item
  Trace the four-part architectural arc of this textbook and map each
  chapter to its corresponding tier in the physical AI systems stack.
\end{enumerate}

\emph{Core chapter takeaway.} The causal boundary is a physical
location, not an abstraction. Below it, the machine shapes its own
future sensory observations, computational delay is paid in physical
displacement, and energy released into the world cannot be recalled.

\end{fbxSimple}

\section{The Industrial Frontier: Artificial Intelligence Beyond
Glass}\label{sec-boundary-1-1}

\lettrine{E}{very}\lettrine{m}{ajor} epoch in computing has been defined
by the boundary across which bits travel. In personal computing, bits
crossed from mainframe backplanes to human eyes through glowing
cathode-ray displays. In cloud computing, bits crossed global
fiber-optic networks to synchronize distributed relational databases. In
modern generative artificial intelligence, bits traverse
trillion-parameter neural networks to generate text, synthesize images,
or recommend financial actions. Yet across all these transformations,
one fundamental property of computer software remained inviolate: the
computation lived safely behind glass.

Today, a massive global industrial race is underway to bring artificial
intelligence beyond the glass. Companies developing humanoid robots
(such as Tesla Optimus, Figure, and Boston Dynamics), autonomous vehicle
fleets (such as Waymo and Zoox), intelligent aerial drones (such as
Skydio), and robotic surgery platforms are deploying deep neural
policies directly aboard mobile, articulated hardware. Yet as these
systems transition from virtual benchmarks to real-world deployment,
engineering teams encounter a fundamental architectural barrier: the
\textbf{Tower of Babel} crisis.

In a typical embodied AI design review, four distinct engineering
cultures clash:

\begin{itemize}
\tightlist
\item
  \textbf{The Machine Learning Perspective:} Focuses on training deep
  transformer backbones and diffusion policies on millions of simulated
  steps, optimizing statistical loss functions and celebrating high
  validation scores, often assuming that real-world deployment is
  primarily a challenge of scaling compute and collecting broader
  teleoperation datasets.
\item
  \textbf{The Systems and Computer Architecture Perspective:} Focuses on
  memory bus bandwidth, DRAM bank conflicts, cache contention, and tail
  latency percentiles, recognizing that streaming multi-gigabyte
  foundation model weights at \(50\text{ Hz}\) threatens to saturate the
  memory bus and starve real-time camera DMA pipelines.
\item
  \textbf{The Robotics and Classical Control Perspective:} Focuses on
  Newton-Euler equations, coordinate transforms, actuator torque
  saturation, and stability margins, recognizing that unconstrained
  neural action chunks risk injecting high-frequency jerk that degrades
  mechanical gearboxes and excites unmodeled resonance.
\item
  \textbf{The Embedded Silicon and Functional Safety Perspective:}
  Focuses on bare-metal interrupt service routines, watchdog timers,
  hardware PWM trip zones, and safety standards (such as ISO 26262 and
  UL 4600), ensuring that non-deterministic neural weights never hold
  unverified direct write access to motor inverter registers.
\end{itemize}

None of these four perspectives can succeed in isolation. You cannot
solve a broken gearbox or an inductive rail voltage collapse with
gradient descent, and you cannot achieve open-vocabulary semantic task
generalization with a hand-tuned PID controller.

\protect\phantomsection\label{callout-definitionux2a-1.2}
\begin{fbxSimple}{callout-definition}{Definition:}{Physical Artificial Intelligence (Physical AI)}
\phantomsection\label{callout-definition*-1.2}

\textbf{Physical Artificial Intelligence (Physical AI)} is the systems
engineering science of embedding high-capacity learned statistical
models into closed-loop physical machines governed by classical
mechanics, thermodynamics, and hard real-time silicon determinism.

\end{fbxSimple}

Physical AI is the systems engineering synthesis that unifies these four
cultures into a single coherent architecture. It trains the
\textbf{Physical AI Systems Architect}---the bilingual systems leader
who designs the complete end-to-end container, ensuring that
high-capacity neural deliberation (Layer 3) safely and deterministically
coordinates with real-time silicon (Layer 2) and physical dynamics
(Layer 1) under formal governance (Layer 4).

To understand why this closed causal loop requires a fundamentally new
systems architecture, we must examine what physically occurs at the
hardware register level when software commands an actuator.

\section{Software That Acts: The Causal
Boundary}\label{sec-boundary-1-2}

Advisory software operates inside an informational container where
errors carry negligible physical momentum. When a recommender model
serves an irrelevant product, a classifier mislabels an image, or a
language model emits an inaccurate token, the computational pipeline
terminates in a display buffer, a network socket, or a sandboxed
database table. The failure remains bounded by the digital medium. A
user dismisses the recommendation, a human analyst closes the alert, or
a database transaction coordinator issues an abort that restores memory
pages from an undo log. \emph{Advisory software is immune to its own
errors because of where it stops, not because of how it was built.} The
safety of the pipeline does not stem from statistical guarantees in the
training objective or rigorous validation on test splits, but from the
physical detachment of the interface. Between the output of the neural
network and any change in the physical world, an air gap or a human
operator intercepts the signal.

Physical artificial intelligence eliminates that detachment. In a
physical machine, the final instruction of a control loop does not write
to a screen or populate a web response. It executes a memory-mapped bus
transaction that updates a hardware register on an actuator controller.
Depending on the machine architecture, that hardware register is a
digital-to-analog converter (DAC) latch holding a commanded reference
voltage, a pulse-width modulation (PWM) compare register (such as
\texttt{TIMx\_CCR} on an ARM Cortex-M or \texttt{ePWM\ CMPA} on a TI
C2000) setting the duty cycle of a power inverter stage, or a serial
control word configuring an H-bridge gate driver.\sidenote{\footnotesize In canonical
  safety microcontrollers, hardware trip zones (\texttt{EPWM\_TZFRC},
  \texttt{TIMx\_BKIN}) operate independently of the CPU. Hardware break
  circuits force complementary PWM outputs into high-impedance or
  low-side clamp states within \(<50\text{ ns}\), preventing
  shoot-through before software ISRs execute.} At that exact memory
address, informational abstraction ends. The bit pattern written to the
register immediately modulates gate voltages across power MOSFETs,
drawing current from a DC bus into the copper windings of a motor.
Current develops magnetic flux, which generates mechanical Lorentz
torque across the rotor. That torque accelerates the rotor and
transmission against reflected load inertia, imparting physical momentum
to the mechanism.

\protect\phantomsection\label{callout-definitionux2a-1.3}
\begin{fbxSimple}{callout-definition}{Definition:}{The Causal Boundary}
\phantomsection\label{callout-definition*-1.3}

The \textbf{causal boundary} is the physical and electrical interface
where digital computational bits convert into physical forces, currents,
or flows. Above the boundary, software state is digital, sandboxed, and
computationally idempotent; below the boundary, actions permanently
alter the physical environment, dissipate irreversible energy, and
endogenously reshape all subsequent sensory observations.

\end{fbxSimple}

When a digital neural network misclassifies an image in a cloud
benchmark, the consequence is confined to an incorrect integer stored in
a memory register. When an autoregressive language model generates an
ungrounded hallucination, the failure manifests as superfluous text on a
browser screen. Software engineers can catch runtime errors, discard
invalid outputs, roll back transactions to a clean checkpoint, or rely
on human users to ignore an answer. In computer systems engineering,
this property is known as \emph{idempotency}: the state of the
surrounding physical universe remains completely decoupled from the
intermediate outputs of the software pipeline.\sidenote{\footnotesize \textbf{For the
  Physical Engineer:} \emph{Idempotency} in digital software means an
  operation can be executed multiple times without changing the final
  state (e.g., refreshing a webpage or retrying a database write). In
  physical systems, no operation is idempotent---every motor command
  consumes irreversible energy and alters the physical environment.}
Memory can be cleared, registers zeroed, and processes restarted without
permanent side effects on the environment.

Physical artificial intelligence begins at the exact instant this
isolation breaks. The moment an algorithmic prediction leaves the
digital domain to modulate a pulse-width modulation motor drive, fire a
valve solenoid, or adjust an electrical grid inverter, computation
ceases to be advice and becomes physical force. At this threshold, an
unhandled exception is no longer a benign runtime error; it is a burnt
motor winding, a sheared gearbox, or a collision with an obstacle in the
physical world.

This transition from digital state to physical force destroys the
statistical foundation of classical machine learning. In offline vision,
natural language processing, and tabular prediction, data collection is
exogenous: the world generates a distribution of inputs, and the model
processes them without altering the mechanism that created them. An
incorrect classification on one image has zero effect on the
distribution governing the next image. In a physical system, the
environment evolves in a closed causal loop: the next physical state and
sensor observation depend directly on the action produced by the prior
model inference.

The sensory input stream is endogenous, generated by the machine's own
past actions. When a policy produces an erroneous torque command, the
joint rotates, the platform shifts, and mounted cameras or lidar units
swing into new viewpoints. The machine drives itself into an unfamiliar
region of the state space where training data is sparse, increasing the
probability of subsequent inference errors and compounding the
divergence from nominal operation. This fundamental architectural
divergence is illustrated in Figure~\ref{fig-01-causal-loop},
contrasting the open-loop, idempotent execution of advisory digital
machine learning with the closed causal loop of Physical AI crossing the
actuator register boundary.

\begin{figure}[t]

\centering{

\includegraphics[width=0.95\linewidth,height=\textheight,keepaspectratio]{/Users/VJ/GitHub/MLSysBook-integrate-physical/physical/book//Users/VJ/GitHub/MLSysBook-integrate-physical/physical/book/chapters/01-boundary/figures/fig01_causal_loop.pdf}

}

\caption{\label{fig-01-causal-loop}\textbf{The Closed Causal Loop of
Physical AI versus Open-Loop Digital Computation.} (a) Open-loop
advisory machine learning operates behind glass on exogenous benchmark
distributions, where computational errors are idempotent and bounded by
software exception handling or transactional rollbacks. (b) In Physical
AI, neural deliberation proposals pass through a deterministic nervous
system safety gate before committing to hardware actuator registers.
Once memory-mapped pulse-width modulation compare registers latch,
electrical current converts into Lorentz force, mechanical torque,
kinetic momentum, and Joule heat, permanently mutating environment state
and endogenously shifting all subsequent sensory observations.}

\end{figure}%

Because physical mass possesses inertia and computation takes finite
time, latency in physical AI translates directly into physical
displacement. As shown on the autonomous testbed in
Figure~\ref{fig-01-sensor-rack}, real-world robotic systems rely on
dense multi-modal sensor suites---including spinning LiDAR scanners,
camera arrays, and radar units---that continuously stream high-bandwidth
telemetry to onboard processors. Every stage in this pipeline---from
photodiode exposure and bus serialization (MIPI CSI-2,
GMSL2),\sidenote{\footnotesize Automotive SerDes links (TI FPD-Link III, ADI GMSL2/3)
  serialize multi-gigabit MIPI CSI-2 sensor streams over coaxial cables
  with sub-\(15\text{ }\mu\text{s}\) transfer latency, hardware CRC-32
  verification, and bidirectional backchannel \texttt{FSYNC} triggers
  that align multi-camera exposure start times within sub-microsecond
  jitter.} to neural inference on accelerator silicon and packet
transmission over fieldbuses (CAN-FD, EtherCAT)---consumes precious
milliseconds. During this aggregate delay window, the machine cannot
incorporate new sensory information and continues moving along the
trajectory dictated by its existing momentum.

\protect\phantomsection\label{callout-ruleux2a-1.4}
\begin{fbxSimple}{callout-rule}{Architect’s Rule of Thumb:}{Latency is Paid in Millimeters}
\phantomsection\label{callout-rule*-1.4}

In advisory machine learning, computational delay only frustrates a
user; in physical AI, latency is paid directly in physical displacement.
At velocity \(v\), an unbudgeted delay \(\Delta t\) consumes unguided
motion that no downstream neural capacity can ever recover.

\end{fbxSimple}

As a concrete example, consider a mobile platform traveling at a modest
walking speed of \(2.0\text{ m/s}\). If the perception pipeline, neural
policy inference, and motor controller communication consume a combined
latency of \(150\text{ ms}\), the platform traverses
\((2.0\text{ m/s}) \times (0.15\text{ s}) = 0.30\text{ m}\)
(\(30\text{ cm}\)) before the first motor response can occur. During
those \(30\text{ cm}\), the vehicle is entirely blind to any
environmental change that occurred after the initial sensor exposure. If
an obstacle enters the path \(20\text{ cm}\) ahead during that inference
window, no amount of model capacity or predictive accuracy can alter the
outcome, because the physical displacement has already been accumulated
by the body.

Once digital commands transition into mechanical work, the state change
becomes thermodynamically irreversible. In standard software systems,
erroneous mutations are contained through exception handling, memory
isolation, and atomic rollback protocols. When an unhandled exception
occurs, the runtime unwinds the call stack, reclaims heap allocations,
and restores the machine state to a known clean checkpoint. The physical
world provides no rollback mechanism. When an actuator controller
applies current to a motor, electrical energy transfers from the battery
into kinetic energy in mechanical linkages or thermal energy dissipated
through resistive Joule heating in the stator windings. If an incorrect
joint torque drives a manipulator into a rigid steel fixture, the
kinetic energy converts into plastic deformation of the tool, mechanical
strain on the gearbox teeth, and heat. No software interrupt can pull
that kinetic energy out of the crushed structure or restore sheared gear
teeth to their original geometry.

\begin{figure}[t]

\centering{

\includegraphics[width=0.85\linewidth,height=\textheight,keepaspectratio]{/Users/VJ/GitHub/MLSysBook-integrate-physical/physical/book/chapters/01-boundary/figures/fig01_real_lidar_sensor_rack.jpg}

}

\caption{\label{fig-01-sensor-rack}\textbf{Spinning Multi-Beam LiDAR
Sensor Hardware for Real-Time Perception.} Close-up photograph of the
64-beam spinning optical LiDAR sensor assembly (Velodyne HDL-64E)
mounted aboard an autonomous research testbed. In Physical AI,
high-bandwidth sensor streams (1.33 million points/second) must be
serialized over real-time interfaces (e.g., FPD-Link III, GMSL2,
Automotive Ethernet), synchronized via hardware timestamps, and
processed through neural perception backbones within strict end-to-end
latency budgets, because every millisecond of transport and compute
delay converts directly into unguided physical displacement. (Source:
Steve Jurvetson, CC BY 2.0).}

\end{figure}%

Understanding why Physical AI requires a distinct systems engineering
discipline requires contrasting its core properties against both
advisory digital machine learning and classical control theory. As
summarized in Table~\ref{tbl-01-paradigm-comparison}, the three
paradigms are sorted from pure software isolation on the left, to pure
classical dynamics in the middle, to the unified physical AI synthesis
on the right.

\begin{longtable}[]{@{}
  >{\raggedright\arraybackslash}p{(\linewidth - 6\tabcolsep) * \real{0.2500}}
  >{\raggedright\arraybackslash}p{(\linewidth - 6\tabcolsep) * \real{0.2500}}
  >{\raggedright\arraybackslash}p{(\linewidth - 6\tabcolsep) * \real{0.2500}}
  >{\raggedright\arraybackslash}p{(\linewidth - 6\tabcolsep) * \real{0.2500}}@{}}
\caption{\textbf{Comparison Matrix of Computing Paradigms across the
Causal Boundary.} Systematic comparison of Advisory Digital Machine
Learning (left), Classical Robotics \& Control Theory (middle), and
Physical Artificial Intelligence (right) across operational substrate,
model formulation, feedback dynamics, execution cadence, actuation
authority, failure semantics, and safety certification
primitives.}\label{tbl-01-paradigm-comparison}\tabularnewline
\toprule\noalign{}
\begin{minipage}[b]{\linewidth}\raggedright
Architectural Dimension
\end{minipage} & \begin{minipage}[b]{\linewidth}\raggedright
Digital Machine Learning (Advisory)
\end{minipage} & \begin{minipage}[b]{\linewidth}\raggedright
Classical Robotics \& Control
\end{minipage} & \begin{minipage}[b]{\linewidth}\raggedright
Physical Artificial Intelligence (Closed-Loop)
\end{minipage} \\
\midrule\noalign{}
\endfirsthead
\toprule\noalign{}
\begin{minipage}[b]{\linewidth}\raggedright
Architectural Dimension
\end{minipage} & \begin{minipage}[b]{\linewidth}\raggedright
Digital Machine Learning (Advisory)
\end{minipage} & \begin{minipage}[b]{\linewidth}\raggedright
Classical Robotics \& Control
\end{minipage} & \begin{minipage}[b]{\linewidth}\raggedright
Physical Artificial Intelligence (Closed-Loop)
\end{minipage} \\
\midrule\noalign{}
\endhead
\bottomrule\noalign{}
\endlastfoot
\textbf{Operational Substrate} & Decoupled virtual memory buffers (GPU
VRAM, DRAM, web endpoints) & Continuous analog state space governed by
differential equations & Hybrid: high-dimensional neural representations
mapped to analog physical actuators \\
\textbf{Model Formulation} & High-capacity statistical approximators
(LLMs, ViTs, diffusion policies) & Analytically specified transfer
functions and deterministic state-space matrices & Foundation policies
paired with deterministic runtime safety filters \\
\textbf{Feedback Dynamics} & Exogenous and static; independent and
identically distributed samples & Continuous deterministic state
feedback with Gaussian sensor noise & Endogenous, non-stationary closed
loop; actions shape future observations \\
\textbf{Execution Cadence} & Throughput-optimized batching; soft or
unbounded latency budgets & Hard real-time single-rate loops on
microcontrollers (100--1000 Hz) & Asynchronous multi-rate cadences:
System 2 (1 Hz), System 1.5 (50 Hz), System 1 (1000 Hz) \\
\textbf{Actuation Authority} & Advisory; sandboxed APIs, display
buffers, or human-in-the-loop validation & Hardcoded feedback directly
driving power inverter timers and PWM registers & Decoupled
Proposal--Permission architecture; untrusted proposals gated by MCU
safety referee \\
\textbf{Failure Semantics} & Idempotent; call-stack unwind, database
rollback, or human alert dismissal & Actuator saturation or tracking
error within analytically bounded stability envelope & Irreversible
thermodynamic work, plastic deformation, kinetic collision, or stator
overheating \\
\textbf{Safety Assurance} & Offline statistical loss minimization on
validation splits & Formal mathematical stability proofs (Lyapunov
functions, gain/phase margins) & Layered runtime invariant enforcement,
forward-invariant barrier filters, and fault injection \\
\end{longtable}

\section{Three Machine Archetypes}\label{sec-boundary-1-3}

Once digital predictions cross the causal boundary into physical force,
they interact with physical machines of every conceivable shape and
scale. An autonomous delivery drone navigating turbulent wind gusts
seems to share nothing with a six-axis robotic arm assembling an
automotive transmission, and neither appears related to a containerized
battery management plant regulating liquid coolant. A practitioner or
student might wonder: how can a single systems engineering curriculum
possibly encompass such wildly diverse physical systems without
dissolving into a collection of bespoke, ad-hoc recipes?

The key is to look past cosmetic form factors and examine the underlying
physical conservation law that binds the machine. Beneath the surface
diversity of shapes, sizes, and applications, every embodied AI system
can be systematically categorized into one of three canonical physical
archetypes based on the physical quantity its actuators manipulate
(Figure~\ref{fig-01-three-archetypes}): moving mass across space,
applying force through physical contact, or regulating continuous fluid
and thermal flows. By viewing physical AI through this three-archetype
lens, a systems architect can immediately determine which resource
budget is exhausted first during a software stall or model anomaly, and
design the real-time nervous system accordingly.

\subsection{Class 1: Mass-Dominated Mobility Systems (Moving Through
Space)}\label{class-1-mass-dominated-mobility-systems-moving-through-space}

Class 1 machines move mass across space, where execution is bound first
by momentum conservation and kinetic energy dissipation
(Figure~\ref{fig-01-three-archetypes} a). In autonomous transport
platforms, quadruped robots navigating rugged terrain, mobile
manipulators, and aerial drones, the actuator bridge applies force to
translate or rotate a body with mass \(m\) at velocity \(v\). Because
kinetic energy \(E_k = \frac{1}{2}m v^2\) cannot be recalled once
delivered to the chassis, safety is measured as remaining spatial
clearance. When braking begins at maximum allowable deceleration
\(a_{\max}\), the physical floor for stopping distance follows the
fundamental kinematic relation:

\[d_{\text{stop}} = v \cdot t_{\text{lag}} + \frac{v^2}{2a_{\max}}\]

where \(t_{\text{lag}}\) is the total sense-to-actuation transport and
inference delay. In a mass-moving system, computational latency converts
directly into unguided spatial displacement
(\(v \cdot t_{\text{lag}}\)), while braking distance scales
quadratically with velocity (\(v^2 / 2a_{\max}\)). The first boundary
breached during a software stall or model misclassification is the
physical distance to an external obstacle.

\subsection{Class 2: Contact-Dominated Manipulation Systems (Applying
Contact
Force)}\label{class-2-contact-dominated-manipulation-systems-applying-contact-force}

Class 2 machines apply mechanical force through contact, where the
primary binding limit is actuator torque saturation, current loop
bandwidth, and structural stiffness
(Figure~\ref{fig-01-three-archetypes} b). In articulated manipulators,
precision assembly fixtures, and robot learning workcells executing
contact-rich insertion and grasping, the end-effector interacts directly
with a rigid or compliant workpiece. Unlike a mobile platform operating
in open space, an arm in contact with a surface cannot resolve an
unexpected displacement command by rolling forward. When an actuator
commands position into a rigid boundary, displacement is constrained and
the interaction is governed by mechanical impedance. By Hooke's law, the
contact force generated by a spatial penetration error \(\Delta x\)
against an environment of stiffness \(k\) is
\(F = k \Delta x\).\sidenote{\footnotesize Rigid contact obeys Signorini-Coulomb
  complementarity (\(0 \le \delta \perp F_N \ge 0\)). Real-time stacks
  approximate compliance via Kelvin-Voigt models to prevent force loop
  divergence.}

Because contact stiffness in metal and tooling fixtures is extremely
high (\(k \ge 10^5\text{ N/m}\)), small geometric trajectory deviations
convert into massive contact forces within milliseconds. If the nervous
system fails to clamp motor current within sub-millisecond response
windows, the motor drives the transmission past its mechanical yield
stress. For example, a tooling tip moving at just \(0.20\text{ m/s}\)
that encounters a stiff workcell fixture will accumulate
\(30\text{ mm}\) of uncorrected displacement during a \(150\text{ ms}\)
software stall---generating thousands of Newtons of force that instantly
shears precision gear teeth. The critical floor in contact systems is
not stopping distance, but the time rate of force rise relative to
actuator current loop bandwidth and structural yield stress.

\subsection{Class 3: Flow-Dominated Process and Energy Systems
(Regulating Energy and Fluid
Flows)}\label{class-3-flow-dominated-process-and-energy-systems-regulating-energy-and-fluid-flows}

Beyond moving linkages and rigid contacts, Class 3 machines regulate
continuous fluid, thermal, or electrical power flows---a domain that
serves both as dedicated process plants and as the foundational
electrical and thermal substrate underlying all high-performance mobile
and manipulation systems (Figure~\ref{fig-01-three-archetypes} c). In
grid-scale battery storage facilities, datacenter liquid cooling loops,
and grid-tied power inverters managing renewable generation, actuators
manipulate coolant pumps, flow valves, and high-frequency power
semiconductors rather than mechanical joints. The state of the system is
governed by mass transport, thermal capacitance, and energy balance
rather than kinematic positions. When a coolant valve or pump stalls,
accumulated heat drives rapid temperature excursions toward chemical
decomposition thresholds. Similarly, a grid-tied inverter operating on a
\(60\text{ Hz}\) utility network faces an electrical phase cycle of
\(16.6\text{ ms}\) (half-cycle zero crossing of \(8.33\text{ ms}\)),
where an unhandled scheduling delay creates alternating-current phase
angle errors that drive destructive surge currents through substation
transformers. The deadline in flow architectures is set by thermal
capacitance, transport delay, and grid phase timing, defining physical
budgets where accumulated energy cannot be quickly reversed once
started.

\subsection{The Integrative Capstone: Multi-Class Embodiments and
Humanoid
Robotics}\label{the-integrative-capstone-multi-class-embodiments-and-humanoid-robotics}

A central question for physical AI systems architects is: \emph{Where do
complex general-purpose embodiments, such as humanoid robots, fit into
this tri-archetype framework?}

A humanoid robot is not a disjoint fourth archetype. Rather, it
represents the \textbf{Integrative Capstone Embodiment}---the ultimate
cyber-physical convergence where all three canonical physical classes
operate and contend simultaneously on a single, tightly constrained
mobile chassis:

\begin{enumerate}
\def\labelenumi{\arabic{enumi}.}
\tightlist
\item
  \textbf{Class 1 Dynamics in Locomotion:} Underactuated bipedal
  balancing, dynamic footstep placement, momentum conservation
  (\(m\mathbf{v}\)), and ground contact slip during walking and stair
  climbing require continuous spatial clearance and kinetic energy
  management.
\item
  \textbf{Class 2 Dynamics in Manipulation:} High-DOF bimanual arms and
  dexterous tactile hands interacting with rigid workpieces, opening
  doors, and grasping tools generate instantaneous contact wrenches
  (\(\mathbf{F} = k \Delta \mathbf{x}\)) bounded by gearbox yield stress
  and joint impedance control bandwidth.
\item
  \textbf{Class 3 Dynamics in Actuation \& Power Delivery:} Driving 20
  to 50 high-torque actuators simultaneously from a single onboard DC
  battery pack pushes inverter thermal dissipation (\(I^2 R\)) to its
  limits and creates severe inductive power-rail voltage droop
  (\(L \frac{dI}{dt}\)) during simultaneous leg shock absorption and
  heavy arm lifting.
\end{enumerate}

While a single-class machine (such as a wheeled AMR or a fixed
industrial arm) isolates its physical budget to a single dominant
conservation law, a humanoid robot forces the systems architect to
manage \textbf{cross-class coupled contention}. When a bipedal robot
catches a falling heavy object, the contact wrench (Class 2) shifts its
center of mass, requiring an instantaneous balancing recovery step
(Class 1) that draws peak phase currents threatening to collapse the
electrical bus and overheat the knee actuators (Class 3). Throughout
this book, we will use individual archetypes to isolate specific
architectural principles in depth, while using humanoid robotics as the
capstone testbed where all physical constraints collide.

\begin{figure}[t]

\centering{

\includegraphics[width=1\linewidth,height=\textheight,keepaspectratio]{/Users/VJ/GitHub/MLSysBook-integrate-physical/physical/book/chapters/01-boundary/figures/fig01_three_archetypes.jpg}

}

\caption{\label{fig-01-three-archetypes}\textbf{The Three Canonical
Machine Archetypes of Physical AI.} (a) Mass-dominated mobility systems
(such as the Boston Dynamics Spot quadruped robot navigating rugged
outdoor terrain) bound first by momentum, kinetic energy dissipation,
and spatial clearance. (b) Contact-dominated manipulation systems (such
as the Franka Emika Panda articulated arm) bound first by contact
impedance, joint torque saturation, and structural yield stress. (c)
Flow-dominated process and thermal plants (such as containerized
grid-scale battery storage systems) bound first by lumped thermal
capacity, transport lag, and electrical phase timing. Complex
embodiments like humanoid robots serve as the integrative capstone
uniting all three classes on a single chassis. (Sources: Wikimedia
Commons, CC BY-SA 4.0).}

\end{figure}%

For any physical AI system, engineering the nervous system begins by
determining which physical budget is exhausted first during a failure.
In software systems, resource constraints center on processor cycles,
memory allocation, and network bandwidth. In physical systems, the
systems architect must identify whether spatial clearance, structural
stress, or heat dissipation creates the critical timing deadline. In a
high-speed vehicle or quadruped, spatial clearance fails within
fractions of a second while motor winding temperatures remain safely
below thermal limits. In an assembly press, mechanical stress exceeds
the yield point of structural members within tens of milliseconds long
before the chassis moves perceptibly or heats up. In a chemical reactor
or battery pack, mechanical stress is negligible, but accumulated
thermal energy breaches safe operational envelopes over several seconds.
The shortest time to failure across these three domains defines the
deterministic execution deadline that the nervous system must meet.

Yet understanding that a machine belongs to one of these three physical
archetypes does not automatically mean it requires the full systems
discipline of Physical AI. A classical hydraulic press or a
deterministic PID drone autopilot manipulates mass and contact force,
but operates entirely within classical control theory. To determine when
a machine crosses into Physical AI, we must establish a rigorous
diagnostic scope test.

\section{The Physical AI Scope Test}\label{sec-boundary-1-4}

Before designing or deploying an embodied architecture, a system
architect or practitioner must establish whether a machine actually
requires the full systems engineering discipline of Physical AI. Without
a rigorous scope boundary, engineering teams make two symmetric
mistakes: they over-engineer classical deterministic machines with
bloated neural models where closed-form PID control is provably optimal,
or they dangerously under-engineer stochastic neural models by granting
them unverified motor authority in safety-critical physical loops.

The scope of Physical AI is defined by the simultaneous intersection of
three structural criteria (Figure~\ref{fig-01-scope-venn}):

\begin{enumerate}
\def\labelenumi{\arabic{enumi}.}
\tightlist
\item
  \textbf{Learned Statistical Model (Unspecifiable Policy):} The system
  relies on a policy, perception encoder, or world model acquired from
  data (e.g., neural networks, foundation models, diffusion policies)
  whose transfer function cannot be written in closed-form analytical
  equations and whose worst-case outputs cannot be mathematically
  bounded offline.
\item
  \textbf{Consequential Physical Feedback:} Actuator actions alter the
  material state of the physical environment, endogenously determining
  what the sensors observe next under the governing laws of classical
  mechanics, thermodynamics, or electromagnetism.
\item
  \textbf{Delegated Physical Authority (Direct Actuator Write):}
  Automated software routines write directly to hardware actuator
  registers without synchronous human approval in the control loop,
  operating within the temporal window where delay or error converts
  directly into irreversible physical momentum.
\end{enumerate}

\begin{figure}[t]

\centering{

\includegraphics[width=0.9\linewidth,height=\textheight,keepaspectratio]{/Users/VJ/GitHub/MLSysBook-integrate-physical/physical/book//Users/VJ/GitHub/MLSysBook-integrate-physical/physical/book/chapters/01-boundary/figures/fig01_scope_venn.pdf}

}

\caption{\label{fig-01-scope-venn}\textbf{The Physical AI Scope Test.}
Physical AI exists strictly at the three-way intersection of learned
models (unspecifiable policies), consequential physical feedback
(actions govern future sensory observations under mechanics and
thermodynamics), and delegated physical authority (automated register
writes without synchronous human intervention). Systems satisfying
exactly two criteria define adjacent engineering fields: Advisory
Decision Support (no physical authority), Classical Control Theory
(analytical transfer functions with Lyapunov certificates), and Digital
Autonomous Systems (transactional databases with idempotency and
rollback). The shared autonomy reflex gap marks the operational boundary
where human neuromuscular reflex latency (\(150\text{--}250\text{ ms}\))
delegates transient authority to the learned model.}

\end{figure}%

A closed feedback loop is necessary for physical interaction, but
feedback alone does not place a machine in this domain. Classical
proportional-integral-derivative controllers run at \(1000\text{ Hz}\)
across aerospace and industrial automation, regulating valve positions
and motor torques through tight sensory feedback. Their feedback loops
are continuous and physical, but their control laws are specified
entirely by fixed mathematical formulas whose stability bounds can be
proven analytically from transfer functions. They hold physical
authority and experience physical feedback, but they lack a learned
model. Conversely, high-frequency algorithmic trading engines execute
automated market orders under learned predictive models within
microseconds. They combine learned policies with delegated transactional
authority, but the feedback they receive is purely digital and
informational. When an order executes, it alters an electronic order
book rather than momentum, friction, or thermal energy. An unhandled
stall in algorithmic trading costs financial capital, but it does not
cause an inertia-driven collision with a structural fixture.

A learned model is equally insufficient on its own to define the field.
Modern engineering workflows frequently deploy deep neural networks to
optimize robot trajectories offline or inspect manufactured parts on
assembly lines. An offline neural trajectory optimizer may compute
time-optimal paths for a six-axis robotic manipulator over several
seconds of cluster compute, but if that trajectory is verified by an
engineer or executed through a classical deterministic controller with
hard-coded envelope stops, the learned model holds no delegated physical
authority over the plant. Similarly, a convolutional neural network
monitoring an assembly line may classify surface defects at
\(60\text{ frames/s}\). If the classification output merely logs an
entry in a database or flags a display for a human operator, the model
operates outside the causal loop. A false negative or an inference
timeout produces a corrupted record, but it cannot directly command an
actuator to over-torque a mechanical joint.

Systems that satisfy exactly two criteria demonstrate where the
structural boundary lies. A real-time driver drowsiness monitor uses a
learned computer vision model to track eyelid movement and vehicle lane
position, receiving continuous physical feedback from the vehicle state.
Yet it lacks delegated physical authority because its only output is an
auditory warning chime. If the vision model fails completely during an
inference stall, the speaker emits no sound, but the steering rack and
brake calipers remain under human control. An analog bimetallic
thermostat possesses delegated physical authority over a gas furnace
valve and operates in continuous feedback with the thermal mass of a
building, but its transfer function is a fixed thermodynamic property of
two bonded metals rather than a learned representation. An automated
trading bot possesses both a learned policy and delegated execution
authority, but its actions propagate through network packets and
databases without interacting with the physical laws of mass and motion.

The most demanding edge cases occur in shared autonomy, where physical
authority is divided dynamically between a human operator and a learned
model. In robotic surgery or advanced driver assistance, a human
operator provides continuous steering or tool guidance while a learned
policy injects corrections to filter hand tremors, avoid anatomical
boundaries, or compensate for vehicle slip. These architectures satisfy
both the learned component and physical feedback criteria. Whether they
fall within the scope of Physical AI depends strictly on the temporal
authority granted to the model. If the learned policy can command even
\(2.0\text{ N}\) of force or adjust steering angle without requiring
human confirmation within the \(200\text{ ms}\) threshold of human
reflex latency, the model possesses delegated physical authority. During
that \(200\text{ ms}\) window, an erroneous model output or a delayed
inference packet acts on the physical world before the operator can
intercede.

Evaluating an ambiguous machine requires a three-step decision rubric
that tests each criterion independently
(Table~\ref{tbl-01-scope-matrix}). The first test determines whether the
system contains a policy component whose mapping is learned from data
rather than specified analytically in closed form. If it does not, the
machine belongs to classical control theory. The second test determines
whether actuator commands alter a physical plant such that subsequent
sensor observations are governed by mechanics, thermodynamics, or fluid
dynamics. If they do not, the machine is a digital machine learning
system. The third test determines whether the learned component writes
commands to physical actuators without a human supervisor or an
independent deterministic barrier capable of vetoing the action before
physical momentum accumulates. If it does not, the machine is an offline
decision-support tool. When all three tests pass, the system operates as
Physical AI, and its engineering must be governed by the laws of the
causal boundary.

\begin{longtable}[]{@{}
  >{\raggedright\arraybackslash}p{(\linewidth - 10\tabcolsep) * \real{0.1667}}
  >{\raggedright\arraybackslash}p{(\linewidth - 10\tabcolsep) * \real{0.1667}}
  >{\raggedright\arraybackslash}p{(\linewidth - 10\tabcolsep) * \real{0.1667}}
  >{\raggedright\arraybackslash}p{(\linewidth - 10\tabcolsep) * \real{0.1667}}
  >{\raggedright\arraybackslash}p{(\linewidth - 10\tabcolsep) * \real{0.1667}}
  >{\raggedright\arraybackslash}p{(\linewidth - 10\tabcolsep) * \real{0.1667}}@{}}
\caption{\textbf{The Physical AI Scope Test Decision Matrix.} Systematic
evaluation of representative engineering systems across the three joint
criteria (Learned Model, Consequential Physical Feedback, and Delegated
Physical Authority), establishing the governing engineering discipline,
limiting physical factors, and primary failure modes for each
candidate.}\label{tbl-01-scope-matrix}\tabularnewline
\toprule\noalign{}
\begin{minipage}[b]{\linewidth}\raggedright
System Candidate / Case
\end{minipage} & \begin{minipage}[b]{\linewidth}\raggedright
Criterion 1: Learned Model (Unspecifiable Policy)
\end{minipage} & \begin{minipage}[b]{\linewidth}\raggedright
Criterion 2: Consequential Physical Feedback
\end{minipage} & \begin{minipage}[b]{\linewidth}\raggedright
Criterion 3: Delegated Physical Authority (Direct Actuator Write)
\end{minipage} & \begin{minipage}[b]{\linewidth}\raggedright
Scope Classification \& Discipline
\end{minipage} & \begin{minipage}[b]{\linewidth}\raggedright
Limiting Physics \& Critical Failure Mode
\end{minipage} \\
\midrule\noalign{}
\endfirsthead
\toprule\noalign{}
\begin{minipage}[b]{\linewidth}\raggedright
System Candidate / Case
\end{minipage} & \begin{minipage}[b]{\linewidth}\raggedright
Criterion 1: Learned Model (Unspecifiable Policy)
\end{minipage} & \begin{minipage}[b]{\linewidth}\raggedright
Criterion 2: Consequential Physical Feedback
\end{minipage} & \begin{minipage}[b]{\linewidth}\raggedright
Criterion 3: Delegated Physical Authority (Direct Actuator Write)
\end{minipage} & \begin{minipage}[b]{\linewidth}\raggedright
Scope Classification \& Discipline
\end{minipage} & \begin{minipage}[b]{\linewidth}\raggedright
Limiting Physics \& Critical Failure Mode
\end{minipage} \\
\midrule\noalign{}
\endhead
\bottomrule\noalign{}
\endlastfoot
\textbf{Autonomous Warehouse AMR (Class 1)} & \textbf{Yes} (Learned
navigation policy / costmap) & \textbf{Yes} (Kinetic momentum, chassis
inertia, and wheel traction) & \textbf{Yes} (Direct motor drive packets
without human in loop) & \textbf{Physical AI} & Kinetic impact; traction
slip and odometry loss \\
\textbf{Articulated Contact Manipulator (Class 2)} & \textbf{Yes}
(Diffusion action chunking for assembly) & \textbf{Yes} (Mechanical
contact impedance and reaction forces) & \textbf{Yes} (Direct joint
torque and velocity register writes) & \textbf{Physical AI} & Structural
yield stress exceedance; gearbox tooth shear \\
\textbf{High-Speed Drone Classical Autopilot} & \textbf{No}
(Analytically tuned 1000 Hz PID / cascaded loop) & \textbf{Yes}
(Aerodynamic lift, drag, and gyroscopic dynamics) & \textbf{Yes} (Direct
PWM motor speed control) & \textbf{Classical Control Theory} & Control
margin erosion; rotor thrust saturation \\
\textbf{High-Frequency Market Maker} & \textbf{Yes} (Deep transformer
sequence predictor) & \textbf{No} (Digital order-book state; memory and
network queues) & \textbf{Yes} (Automated limit order execution over 10
GbE) & \textbf{Digital Machine Learning / FinTech} & Financial capital
loss; market flash crash (idempotent rollback) \\
\textbf{Advisory Driver Drowsiness Monitor} & \textbf{Yes} (Facial
landmark vision network) & \textbf{Yes} (Vehicle kinematics and driver
physiology) & \textbf{No} (Auditory warning buzzer and dashboard LED
only) & \textbf{Advisory Decision Support} & Human perception-reaction
latency; zero physical authority \\
\textbf{Offline Trajectory Optimizer} & \textbf{Yes} (Deep neural
network for time-optimal path search) & \textbf{No} (Static CAD workcell
model; offline batch execution) & \textbf{No} (Motion plan verified and
loaded via human engineer) & \textbf{Offline Optimization Tool} &
Numerical non-convergence; zero direct runtime plant coupling \\
\textbf{Shared-Autonomy Surgical Assist} & \textbf{Yes} (Learned haptic
tremor filter / virtual fixtures) & \textbf{Yes} (Anatomical tissue
compliance and reaction forces) & \textbf{Conditional} (Model injects
forces within human reflex window) & \textbf{Borderline / Shared
Physical AI} & Human neuromuscular reflex latency vs.~tool force rise
rate \\
\end{longtable}

\section{The Four Pillars of an Embodied Agent}\label{sec-boundary-1-5}

Before exploring how physical systems fail or examining the silicon
circuits that control them, we must map the complete anatomical
landscape of an embodied machine. Every physical AI system---whether a
humanoid biped, an autonomous tractor, or an automated surgical
console---is composed of four fundamental architectural pillars that
operate across distinct physical timescales:

\begin{enumerate}
\def\labelenumi{\arabic{enumi}.}
\tightlist
\item
  \textbf{Pillar 1: The Dynamical Body (Layer 1):} The physical plant
  itself. Comprising mechanical linkages, structural mass, gear
  transmissions, motor windings, power electronic inverters, and battery
  packs. The body does not obey software abstractions, but rather the
  governing conservation laws of classical mechanics, thermodynamics,
  and electromagnetism. It accumulates fatigue, dissipates resistive
  Joule heat (\(I^2 R\)), and resists instantaneous acceleration through
  reflected rotor inertia (\(N^2 J\)).
\item
  \textbf{Pillar 2: The Cognitive Brain (Layer 3):} The high-capacity
  learned intelligence. Comprising multi-modal vision-language-action
  (VLA) foundation models, diffusion policies, and semantic world models
  running on high-performance compute accelerators (GPUs, NPUs). The
  brain provides open-vocabulary generalization, spatial scene
  understanding, and long-horizon affordance reasoning. Its competence
  arises from high-dimensional statistical approximation across massive
  datasets, meaning it is fundamentally stochastic, non-deterministic,
  and prone to out-of-distribution hallucinations.
\item
  \textbf{Pillar 3: The Real-Time Nervous System (Layer 2):} The
  deterministic cyber-physical bridge. Comprising hard real-time
  microcontrollers (bare-metal MCUs), deterministic fieldbus networks
  (CAN-FD, EtherCAT, TSN Ethernet), and lock-free shared memory fabrics.
  The nervous system enforces sub-millisecond execution deadlines
  (\(1000\text{ Hz}\)), guarantees microsecond clock synchronization,
  evaluates candidate commands, and retains exclusive physical authority
  over actuator registers.
\item
  \textbf{Pillar 4: The Governance Envelope (Layer 4):} The formal
  safety, verification, and certification harness. Comprising
  forward-invariant safety monitors (Control Barrier Functions),
  decaying cryptographic operating leases, cross-layer hardware fault
  injection, and auditable Claim-Argument-Evidence (CAE) safety cases
  that qualify the complete machine for deployment.
\end{enumerate}

These four pillars form an integrated physical container. The Cognitive
Brain (Pillar 2) generates high-level intentions and proposed action
trajectories; the Real-Time Nervous System (Pillar 3) evaluates those
proposals against the physical constraints of the Dynamical Body (Pillar
1); and the Governance Envelope (Pillar 4) bounds the operational
envelope and provides formal assurance.

When any one of these four pillars is missing or compromised, the causal
boundary breaks. To see this catastrophic failure mode in forensic
detail, we examine one of the most famous real-world failures in
autonomous systems history.

\section{Anatomy of a Crash: When the Boundary
Fails}\label{sec-boundary-1-6}

To understand why the boundary conditions of Physical AI are not mere
theoretical classifications, but life-critical systems requirements, we
must examine what happens when an embodied machine violates them. When a
high-capacity learned model is granted unverified motor authority over
physical mass, software failure does not remain behind glass---it
converts into kinetic destruction. A physical crash in an autonomous
machine is rarely caused by an inaccurate prediction alone, but by the
absence of an architectural mechanism with the authority to refuse one.

\subsection{The 2018 Tempe Autonomous
Collision}\label{the-2018-tempe-autonomous-collision}

On a clear night in March 2018, an automated test vehicle---a modified
Volvo XC90 SUV operating Uber Advanced Technologies Group's
developmental Automated Driving System (ADS)---was traveling at
forty-three miles per hour (\(v_0 = 19.2\text{ m/s}\), mass
\(m \approx 2100\text{ kg}\)) along a public multi-lane roadway in
Tempe, Arizona (\citeproc{ref-ntsb2019uber}{National Transportation
Safety Board 2019})\sidenote{\footnotesize NTSB Report HAR-19/03
  (\citeproc{ref-ntsb2019uber}{National Transportation Safety Board
  2019}) revealed that disabling the factory autonomous emergency
  braking system eliminated the deterministic physical safety enforcer,
  relying on unmonitored human supervision.}. With a total kinetic
energy of:

\[E_k = \frac{1}{2} m v_0^2 = \frac{1}{2}(2100\text{ kg})(19.2\text{ m/s})^2 \approx 387\text{ kJ}\]

the moving vehicle struck and fatally injured a pedestrian pushing a
bicycle across the asphalt.

The vehicle was equipped with an expansive, state-of-the-art physical
sensor suite (Figure~\ref{fig-01-tempe-vehicle}): a roof-mounted 64-beam
Velodyne LiDAR (360° azimuth, \(> 100\text{ m}\) range), ten optical
cameras, and eight medium- and long-range radar units. At timestamp
\(T - 5.6\text{ s}\), while the vehicle was still \(82.0\text{ m}\)
away, the primary LiDAR and radar sensors established positive contact
with the obstacle, initializing a dynamic track in host computer memory.

\subsection{Perception Toggling and State
Erasure}\label{perception-toggling-and-state-erasure}

Despite continuous sensor contact over more than five seconds, the
vehicle failed to brake. Between \(T - 4.4\text{ s}\) and
\(T - 1.2\text{ s}\), the stochastic neural perception classifier
repeatedly toggled its semantic classification label between an
\emph{unknown vehicle}, a \emph{bicycle}, and an \emph{unclassified
static obstacle}.

Because the software tracking architecture associated Kalman filter
state estimators with discrete semantic categories, each classification
label switch caused the runtime to discard the existing kinematic track
and re-initialize a new filter instance with an assumed velocity vector
of zero. As a direct consequence of this classification instability, the
downstream trajectory planner never accumulated persistent evidence of
an imminent collision course.

\subsection{Latency Suppression and Kinematic
Depletion}\label{latency-suppression-and-kinematic-depletion}

At \(T - 1.2\text{ s}\), the software finally consolidated the track but
triggered an internal \(1.2\text{ s}\) ``suppression delay''---a
software counter designed by developers to prevent false-positive
emergency braking alarms from disrupting passenger comfort.

During this suppression window, the vehicle traversed an unguided
displacement of:

\[d_{\text{lag}} = v_0 \cdot t_{\text{suppress}} = (19.2\text{ m/s}) \times (1.2\text{ s}) \approx 23.0\text{ m}\]

When the planner finally completed the suppression window and requested
emergency braking at \(T - 0.2\text{ s}\), the remaining physical
clearance had collapsed to just \(3.8\text{ m}\). Bringing a
\(2100\text{ kg}\) vehicle to a complete halt on dry asphalt (assuming a
maximum deceleration \(a_{\max} \approx 8.0\text{ m/s}^2\)) requires a
physical braking distance of:

\[d_{\text{brake}} = \frac{v_0^2}{2 a_{\max}} = \frac{(19.2\text{ m/s})^2}{2 \times 8.0\text{ m/s}^2} \approx 23.0\text{ m}\]

Total stopping distance required at \(T - 1.2\text{ s}\) exceeded
\(d_{\text{stop}} = d_{\text{lag}} + d_{\text{brake}} = 23.0\text{ m} + 23.0\text{ m} = 46.0\text{ m}\).
Because available clearance had eroded to \(3.8\text{ m}\), kinetic
collision at full cruising speed was mathematically sealed two full
seconds before impact.

\begin{figure}[t]

\centering{

\includegraphics[width=0.95\linewidth,height=\textheight,keepaspectratio]{/Users/VJ/GitHub/MLSysBook-integrate-physical/physical/book/chapters/01-boundary/figures/fig01_real_ntsb_tempe_vehicle.jpg}

}

\caption{\label{fig-01-tempe-vehicle}\textbf{Autonomous Test Vehicle
Sensor Configuration from NTSB Collision Investigation.} The modified
2017 Volvo XC90 sport utility vehicle operating Uber ATG's developmental
Automated Driving System (ADS) during the March 2018 fatal collision in
Tempe, Arizona. The physical sensor suite comprised a roof-mounted
64-beam Velodyne LiDAR (360° azimuth, \textgreater100 m range), a
10-camera optical sensing suite including forward-facing stereo and
wide-angle cameras, and eight medium- and long-range radar units
providing 360° perimeter coverage. Despite continuous sensor contact
starting at 82 meters (\(T-5.6\text{ s}\)), classification instability
in the perception pipeline and an unverified 1.2-second suppression
delay prevented the planner from requesting braking until 0.2 seconds
before impact. (Source: NTSB Highway Accident Report NTSB/HAR-19/03
(\citeproc{ref-ntsb2019uber}{National Transportation Safety Board
2019}), Public Domain).}

\end{figure}%

\protect\phantomsection\label{callout-autopsyux2a-1.5}
\begin{fbxSimple}{callout-autopsy}{Incident Autopsy:}{The Missing Safety Referee}
\phantomsection\label{callout-autopsy*-1.5}

\begin{itemize}
\tightlist
\item
  \textbf{System:} Developmental Automated Driving System (ADS) on a
  2017 Volvo XC90 SUV (\(2100\text{ kg}\), \(19.2\text{ m/s}\)) equipped
  with roof-mounted 64-beam LiDAR, radar, and optical cameras.
\item
  \textbf{Incident:} Struck and fatally injured a pedestrian pushing a
  bicycle across an urban roadway without executing emergency braking
  before impact.
\item
  \textbf{Telemetry Snapshot:} Sensor contact established at
  \(T - 5.6\text{ s}\) (\(82\text{ m}\)); classifier toggled across
  three semantic labels for \(3.2\text{ s}\); \(1.2\text{ s}\)
  false-alarm suppression delay expired at \(T - 0.2\text{ s}\)
  (\(3.8\text{ m}\) remaining clearance).
\item
  \textbf{Root Cause:} Semantic classification toggling caused the
  tracker to repeatedly discard kinematic state histories and
  re-initialize velocity estimators to zero, preventing the trajectory
  planner from recognizing an imminent collision.
\item
  \textbf{Architectural Flaw:} Complete absence of an isolated
  deterministic safety referee running on dedicated real-time silicon
  with the authority to evaluate geometric clearance invariants directly
  against raw sensor points.
\item
  \textbf{Provenance:} National Transportation Safety Board Accident
  Report NTSB/HAR-19/03 (2019).
\end{itemize}

\end{fbxSimple}

This investigation autopsy demonstrates why refining perception models
or expanding offline training datasets cannot compensate for an
architectural absence of authority boundaries. An open-world statistical
model will inevitably encounter out-of-distribution scenes, sensor
degradations, and transient classification oscillations. The Tempe
vehicle crashed not because its neural network was imperfect, but
because the autonomy architecture lacked an isolated nervous system
running on dedicated real-time silicon with the independent authority to
evaluate spatial clearance and execute an emergency stop. If a
deterministic safety enforcer had continuously checked the stopping
clearance invariant at the hardware floor, the machine would have
executed emergency braking at \(T - 2.5\text{ s}\), bringing the vehicle
to a controlled halt long before available clearance
collapsed.\sidenote{\footnotesize ISO 26262-9 Clause 5 defines Freedom From
  Interference (FFI) across ASIL boundaries. Proving FFI requires
  verified spatial partitioning (MPU/IOMMU memory protection), temporal
  partitioning (hardware bandwidth limiters and watchdog timers), and
  authenticated communication contracts.}

The forensic breakdown of the Tempe crash reveals four universal bedrock
laws that govern all physical machines across the causal boundary.

\section{The Four Bedrock Laws}\label{sec-boundary-1-7}

Four bedrock laws govern every machine whose nervous system translates
compute cycles into physical action. Notice how these four laws map
symmetrically to the \textbf{Four Pillars of an Embodied Agent} and the
\textbf{Four Parts of this Textbook}:

\protect\phantomsection\label{callout-lawux2a-1.6}
\begin{fbxSimple}{callout-law}{Bedrock Law:}{1. Physical Causality and Irreversibility (The Body Law — Part I)}
\phantomsection\label{callout-law*-1.6}

Software state can be caught with an exception block, isolated in a
virtual container, or restored from a checkpoint. Physical actions
governed by mass, inertia, friction, and gravity are permanent and
irreversible. Every actuation permanently alters the physical
environment, dissipates energy, and alters all future sensory
observations endogenously.

\end{fbxSimple}

The first bedrock law is the thermodynamic irreversibility of physical
action. In a pure software system, an erroneous operation is reversed by
rolling back a database transaction, restoring a file from a snapshot,
or restarting a process with clean memory. In a physical machine, an
actuator converts electrical charge into kinetic energy, structural
strain, or fluid pressure. Once that energy crosses the causal boundary
into the environment, it cannot be recalled. When an erroneous command
accelerates a warehouse robot toward an obstacle, the kinetic energy
deposited into the chassis can only be counteracted by applying opposing
forces through friction brakes, which dissipates the energy as heat. If
an over-current command drives a six-axis manipulator into an assembly
jig, the mechanical stress exceeds the yield point of the metal within
milliseconds, permanently deforming the tooling. Software transactions
can be rolled back, but thermodynamic work cannot be undone.

\protect\phantomsection\label{callout-lawux2a-1.7}
\begin{fbxSimple}{callout-law}{Bedrock Law:}{2. Endogenous Data and Compounding Drift (The Policy Law — Part II)}
\phantomsection\label{callout-law*-1.7}

A physical agent shapes its own future data distribution. Naive
behavioral cloning on open-loop benchmarks suffers compounding covariate
drift (\(\mathcal{O}(T^2 \epsilon)\)) in closed-loop execution. An error
on step \(t\) creates an out-of-distribution observation at step
\(t+1\), causing policy confidence to collapse into mechanical
divergence.

\end{fbxSimple}

The second bedrock law governs how embodied policies learn and fail. In
traditional computer vision or natural language processing, data is
exogenous: whether a model correctly identifies an image has zero effect
on the distribution of subsequent images. In Physical AI, every torque
command rotates joints, shifts camera angles, and alters contact forces.
The data stream is endogenous---created entirely by the machine's own
past actions. A single erroneous inference drives the hardware into an
unvisited state where neural predictions exhibit higher variance,
compounding errors until the machine exceeds physical stability limits.

\protect\phantomsection\label{callout-lawux2a-1.8}
\begin{fbxSimple}{callout-law}{Bedrock Law:}{3. Proposal–Permission Privilege and Freshness (The Nervous System Law — Part III)}
\phantomsection\label{callout-law*-1.8}

No stochastic neural network may hold direct motor authority. Learned
models running on application processors are unprivileged proposal
services. Real-time physical permission belongs exclusively to dedicated
microcontrollers executing deterministic safety filters at 1000 Hz.
Latency is physical displacement; stale proposals older than their
freshness deadline must be discarded.

\end{fbxSimple}

The third bedrock law requires strict separation between proposal and
permission. The subsystem of the nervous system that generates complex,
goal-directed behavior operates by optimization over high-dimensional
learned representations. Because deep neural networks are
uninterpretable and produce erratic outputs under out-of-distribution
inputs, a learned policy must never hold direct write access to physical
actuator registers. The role of the learned model is strictly to propose
candidate trajectories. The authority to grant permission, clip commands
to safe physical envelopes, or reject the proposed action entirely must
belong to a separate, deterministic supervisory layer running on
isolated execution hardware.

\protect\phantomsection\label{callout-lawux2a-1.9}
\begin{fbxSimple}{callout-law}{Bedrock Law:}{4. Governed Authority and Defensible Release (The Governance Law — Part IV)}
\phantomsection\label{callout-law*-1.9}

Physical AI systems are never `generally safe'; operating authority is a
decaying lease strictly conditioned on bounded evidence. A cherry-picked
video demonstration is not evidence of safety; defensible release
requires cross-layer hardware fault injection, formal barrier
guarantees, and an auditable safety case.

\end{fbxSimple}

The fourth bedrock law dictates that deployment readiness cannot be
justified by empirical validation loss. In cloud computing, software
reliability is evaluated statistically: an engineer measures mean
accuracy, accepts an uptime target of \(99.9\%\), and resolves rare
corner-case crashes through automated restarts and rolling updates. In
physical AI, the tail of the error distribution manifests as mechanical
destruction. Safe operation cannot be certified by average-case
performance metrics gathered over finite training sets. It requires
formal, worst-case envelope bounds verified against the physical laws of
the plant, ensuring that even if the learned proposal violates safety
limits, the machine remains within guaranteed physical boundaries.

The four bedrock laws establish the physical constraints that govern
every machine. But of these four, the Law of Proposal--Permission
Privilege dictates how we must architect and physically build the
computer system. How do we take an untrusted, stochastic neural policy
and safely translate its high-level intelligence into physical motion on
real hardware?

\section{The Dual-Brain Architecture}\label{sec-boundary-1-8}

To safely bridge high-capacity learned policies and real-time physical
machines, the systems architect implements the
\textbf{Proposal--Permission Architecture} across heterogeneous silicon.
The system partitions compute into two physically separated execution
domains: an unprivileged application processor (Linux MPU or Edge NPU)
generating candidate action proposals at deliberative cadences
(\(10\text{--}50\text{ Hz}\)), paired with an isolated, zero-allocation
real-time microcontroller (bare-metal MCU) executing deterministic
safety filters and physical fallbacks at \(1000\text{ Hz}\) with
exclusive electrical authority over actuator drive stages.

Authority over actuation must belong to a mechanism that can refuse
execution even when the proposer is broken. This separation between
proposal and permission is a foundational structural principle: in a
pure software system, an unprivileged process can be granted direct
write access to application memory because an erroneous write costs only
compute cycles and memory bandwidth to correct. In a physical machine,
state changes move mass, dissipate heat, and expend momentum that cannot
be recalled. Because a learned policy is an inductive statistical model
whose worst-case outputs under distribution shift cannot be
mathematically bounded, the neural network must function strictly as an
unprivileged proposer. Routing learned model outputs directly to
actuator registers couples the mechanical integrity of the machine to
the unverified tail of a probability distribution.

Training-time optimization cannot substitute for runtime invariant
enforcement. In reinforcement learning and imitation learning pipelines,
developers frequently attempt to enforce safety by augmenting the loss
function with penalty terms, barrier losses, or negative reward shaping.
These numerical penalties influence gradient descent during training,
reducing the empirical frequency of undesirable states across the
training distribution. But they provide zero hard guarantees at runtime.
When the deployed policy encounters novel sensor noise, lens occlusion,
or out-of-distribution state combinations, offline optimization cannot
prevent the network from generating an unsafe control vector. Soft
optimization incentives shape probability distributions, but physical
safety requires deterministic boundary enforcement. Safety constraints
must therefore live downstream of the model as an explicit, hard
architectural gate positioned between the policy output and the actuator
registers.

In computer systems engineering, developers often ask: \emph{why can't
we run this proposal--permission gate inside the same high-performance
Linux processor that runs the neural network?} The answer lies in the
fundamental conflict between high-capacity statistical computing and
hard real-time determinism. A modern vision-language-action model or
diffusion policy requires billions of floating-point operations and
gigabytes of memory bandwidth, necessitating multi-core application
processors (MPUs) or neural accelerators (NPUs) running complex
operating systems like Linux. However, general-purpose operating systems
introduce non-deterministic scheduling preemption, dynamic memory
allocation locks, and interrupt latencies ranging from \(10\text{ ms}\)
to over \(100\text{ ms}\). In a physical system, a \(50\text{ ms}\) OS
scheduling hiccup is not a dropped animation frame; it is
\(50\text{ ms}\) of unguided kinetic displacement that drives mechanical
linkages past structural limits. Conversely, one cannot run
high-capacity neural models inside a microcontroller (MCU), because
microcontrollers lack the memory bandwidth and tensor execution units
required for high-dimensional inference.

Therefore, every robust physical AI machine physically partitions its
compute across heterogeneous silicon (Figure~\ref{fig-01-dual-brain}):

\begin{enumerate}
\def\labelenumi{\arabic{enumi}.}
\tightlist
\item
  \textbf{The Cognitive Proposer (Linux MPU / Edge NPU):} Operates in
  unprivileged user space at a deliberative cadence
  (\(10\text{--}50\text{ Hz}\)). It executes high-capacity neural
  policies, multi-modal perception backbones, and trajectory generators.
  It possesses \textbf{zero direct electrical wiring to physical
  actuator registers}. It can only write candidate trajectory proposals
  into a shared SRAM mailbox or lock-free circular ring buffer across an
  inter-processor communication (IPC) channel.
\item
  \textbf{The Real-Time Permission Referee (Bare-Metal / RTOS MCU):}
  Operates on isolated, dedicated silicon at a hard real-time cadence
  (\(1000\text{ Hz}\), \(1\text{ ms}\) deadline,
  \(< 50\text{ }\mu\text{s}\) jitter) with static, zero-allocation
  memory (\texttt{malloc} is strictly forbidden at runtime). It holds
  \textbf{exclusive physical authority} over the PWM compare registers,
  digital-to-analog converters, and gate drivers. The MCU continuously
  checks physical safety invariants and stopping clearance. If the
  proposed command satisfies all physical constraints, the MCU latches
  the values to hardware; if the proposal is unsafe, stalls, or panics,
  the MCU vetoes the command and executes a deterministic physical
  fallback.
\end{enumerate}

To bridge the asynchronous Linux OS and deterministic bare-metal
firmware without priority inversion or lock contention, the MPU and MCU
communicate across a shared SRAM mailbox using lock-free atomic sequence
locks (\emph{seqlocks}) or double-buffered ping-pong ring
buffers.\sidenote{\footnotesize In heterogeneous SoCs (e.g., TI Sitara AM64x, NXP
  i.MX8M), Linux MPUs and real-time Cortex-M/R MCUs exchange data across
  non-cached dual-port SRAM via OpenAMP/RPMsg mailboxes. Out-of-order
  MPU cores require explicit Data Memory Barriers (\texttt{DMB}) and
  cache flushes before asserting Inter-Processor Interrupts
  (\texttt{IPI}) to prevent torn memory reads.} The MPU publishes
candidate action chunks with a monotonically incrementing sequence
version; the MCU reads the trajectory atomically without ever blocking
on a mutex. If an MPU core stalls or experiences an OS scheduling delay,
the sequence version fails to update, immediately signaling the MCU that
proposed commands have aged beyond their freshness deadline.

\begin{figure}[t]

\centering{

\includegraphics[width=1\linewidth,height=\textheight,keepaspectratio]{/Users/VJ/GitHub/MLSysBook-integrate-physical/physical/book//Users/VJ/GitHub/MLSysBook-integrate-physical/physical/book/chapters/01-boundary/figures/fig01_dual_brain_soc.pdf}

}

\caption{\label{fig-01-dual-brain}\textbf{The Physical AI Dual-Brain
Privilege Separation Architecture.} Heterogeneous silicon partitioning
between the untrusted proposal engine (Linux MPU / Edge NPU) executing
high-capacity foundation models and the trusted permission authority
(bare-metal real-time MCU) executing zero-allocation 1 kHz safety
filters over shared SRAM mailboxes. The permission authority exercises
sole veto over actuator phase current setpoints, guaranteeing physical
invariants even during total model stalls or numerical divergence.}

\end{figure}%

\subsection{The Three Real-Time Referee
Gates}\label{the-three-real-time-referee-gates}

Because a machine learning model running on an application processor is
strictly a proposal generator, exclusive authority to commit physical
register writes belongs to the real-time referee. The referee operates
on dedicated microcontroller silicon at a fixed \(1\text{ kHz}\)
interrupt cadence (\(1.0\text{ ms}\) deadline) and enforces three
mandatory gates before latching any command to motor drivers:

\begin{enumerate}
\def\labelenumi{\arabic{enumi}.}
\tightlist
\item
  \textbf{Temporal Freshness Gate:} If the proposal timestamp is older
  than the strict freshness deadline
  (\(\Delta t > \Delta t_{\text{fresh}}\)), the referee drops the stale
  proposal.
\item
  \textbf{Physical Invariant Gate:} If the commanded torque, velocity,
  or trajectory breaches dynamic stopping clearance
  (\(d_{\text{stop}}\)) or current limits, the referee mathematically
  projects the candidate command back to the safe admissible boundary.
  Because microcontrollers have strict cycle budgets, this projection is
  computed using lightweight, closed-form Control Barrier Functions
  (CBFs) or box-constrained quadratic programs (QP) solved in
  microseconds via static, zero-allocation active-set routines without
  non-deterministic matrix inversions.
\item
  \textbf{Deterministic Fallback Gate:} If the neural proposer crashes,
  hangs, or outputs numerical anomalies (\(\text{NaNs} / \pm\infty\)),
  the referee autonomously engages verified regenerative braking or
  passive dynamic centering.
\end{enumerate}

When a candidate action is refused, the permission authority must
execute a deterministic intervention strategy that maintains the
physical machine within its safe operating envelope. Refusal cannot
consist of simply dropping the command and allowing actuator registers
to float. If a mobile base moving at \(2.0\text{ m/s}\) suddenly
receives zero motor current without controlled deceleration, momentum
drives the unconstrained mass into obstacles or induces wheel slip that
invalidates odometry. The permission authority applies a graded
hierarchy of interventions matched to the severity of the violation:

\begin{itemize}
\tightlist
\item
  \emph{Rate Clamping \& Admissible Projection:} For minor envelope
  exceedances, the gate mathematically projects the candidate vector to
  the nearest admissible boundary point, enforcing current limits and
  voltage ceilings.
\item
  \emph{Kinematic Slicing:} For dynamic trajectory proposals, the gate
  clips acceleration and jerk to maximum allowable mechanical
  transmission bounds.
\item
  \emph{Deterministic Physical Fallback:} When the proposer suffers a
  numerical exception (NaNs/Infinities), latency timeout, or heartbeat
  loss, the MCU engages an autonomous physical fallback---commanding
  controlled regenerative braking or executing a verified deceleration
  trajectory to bring the mass to a stationary, passive state.
\end{itemize}

Understanding this dual-brain realization of proposal and permission is
essential because it fundamentally restructures how we develop, verify,
and govern the entire machine. In traditional software services,
engineering teams treat functional safety, cyber-security, and audit
logging as three completely isolated disciplines handled by separate
organizations across disjoint software layers. In Physical AI, the
dual-brain proposal--permission boundary collapses all three concerns
onto a single physical primitive.

\section{Safety Convergence}\label{sec-boundary-1-9}

Safety asks whether the machine caused harm. Security asks whether an
external agent forced it to. Accountability asks whether an auditor can
determine afterwards what occurred and why. In software services, these
three concerns belong to separate engineering organizations and distinct
runtimes, dividing along application logic, network firewalls, and audit
logging databases. In a machine that moves mass and commands current,
all three questions collapse into a single physical primitive at the
causal boundary, determining whether an independent gate possessed the
state and the authority to refuse a register write before torque reached
the plant.

Safety failures arise without an adversary. A camera lens accumulates
dust over \(500\text{ h}\) of operation, shifting the input distribution
away from the training domain until an attention head generates an
out-of-distribution steering torque. A joint bearing develops mechanical
play, introducing unmodeled compliance that converts a stable learned
position trajectory into a resonant limit cycle. Or a sudden change in
surface friction drops tire traction below the friction coefficient
assumed by the learned policy. In each case, the proposer generates an
action that is statistically plausible to the network but physically
hazardous to the body. If the permission gate lacks the independent
kinematic model or thermal budget to detect the exceedance, the
unrefused command writes to the motor drive registers, and stored
kinetic energy is released into the workspace.

Security failures introduce an adversary who deliberately exploits the
gap between statistical inference and physical reality. An attacker does
not need to compromise cryptographic keys or gain root access on the
host operating system if they can manipulate sensory inputs. Projected
structured light can spoof a depth sensor, acoustic injection can
resonate MEMS gyroscopes at their resonant frequency of
\(27\text{ kHz}\) to corrupt inertial odometry,\sidenote{\footnotesize Resonant
  acoustic frequencies (\(18\text{--}32\text{ kHz}\)) can saturate MEMS
  gyroscopes, requiring physical damping and MCU kinematic cross-checks.}
or adversarial visual perturbations on a factory floor can induce a
vision-language-action model to output maximum actuator velocity. On a
traditional server, an injection attack compromises data confidentiality
or integrity within an abstract memory space. In physical AI, an
injection attack commands physical work, turning an unvalidated
inference into mechanical force. The security perimeter cannot terminate
at the network interface or the operating system system-call boundary.
It terminates at the permission gate, which must evaluate every
candidate command against physical invariants regardless of whether the
proposer acted in good faith or under adversarial control.

Traditional software authorization models fail in physical machines
because capability tokens and access control lists assume discrete,
memoryless transactions. In a web service, possessing a valid credential
grants permission to write a record to a database, and the cost of the
operation is bounded by compute cycles and storage bytes. In physical
space, permission is not a static binary entitlement. In an analytical
robotic joint, a command to apply \(50\text{ N}\cdot\text{m}\) of torque
is valid when the link is stationary at mid-stroke, but the identical
command causes structural yield if the joint is \(2\text{ mm}\) from its
hard stop moving at \(1.5\text{ rad/s}\). Physical authority is
state-dependent, continuous, and dynamic, governed by velocity,
momentum, thermal accumulation, and spatial occupancy. A permission gate
cannot evaluate an action in isolation as a stateless remote procedure
call. It must integrate the action over time against the current
physical state of the body, validating that the applied forces remain
within the safe invariant reachable set of the plant.

Accountability requires proving the exact causal chain that led to
actuation when a machine damages hardware or breaches an envelope. In
standard computing systems, logging is asynchronous, best-effort, and
coarse-grained, dropping records during buffer pressure or recording
timestamps with millisecond jitter. At the causal boundary,
accountability demands deterministic, high-rate telemetry captured
synchronously with the control loop. If a joint actuator commands
\(1000\text{ Hz}\) updates with a \(1\text{ ms}\) deadline, the
permission gate must record the raw candidate vector proposed by the
model, the state estimate used by the gate, the mathematical criterion
applied by the safety filter, and the exact vector written to the
hardware registers. Sealing these records in tamper-evident, append-only
circular buffers in non-volatile SRAM/FRAM provides post-incident
reconstruction. When an anomaly occurs, an engineer can verify
mathematically whether the failure originated in a corrupted sensor
reading, a degraded policy proposal, an erroneous permission rule, or a
mechanical component that failed to produce the commanded
counter-torque.

\section{Architecture of This Book: The Journey
Ahead}\label{sec-boundary-1-10}

Engineering dependable physical AI systems requires mastering how
statistical learning, physical dynamics, and hard real-time computer
systems compose across every tier of the machine. To build this
discipline systematically from first principles, this textbook is
organized as an unbroken intellectual journey across four foundational
movements (Table~\ref{tbl-01-book-roadmap}):

\subsection{The Four Movements of Physical
AI}\label{the-four-movements-of-physical-ai}

\textbf{Part I: The Physical Container (Chapter~\ref{sec-boundary}
through Chapter~\ref{sec-nervous}).} Before writing a single line of
policy code or launching a neural network training cluster, an engineer
must first step onto the hardware floor. We begin by characterizing what
the physical world will not negotiate---the unyielding laws of momentum,
reflected gear inertia, thermal dissipation, and microsecond clock
jitter that bound the body, the brain, and the real-time nervous system.
By treating the physical plant as an immutable container, we establish
the non-negotiable temporal and dynamic budgets that every upstream
algorithm must satisfy.

\textbf{Part II: The Policy Lifecycle (Chapter~\ref{sec-data} through
Chapter~\ref{sec-evaluation}).} Once the physical container is defined,
we confront where an embodied policy's intelligence comes from. We trace
the messy, endogenous reality of physical data collection, evaluate why
naive behavioral cloning suffers compounding drift in closed-loop
deployment, and establish why statistical benchmark accuracy is
powerless to guarantee physical survival without rigorous,
non-asymptotic testing and an understanding of the Butler \& Finelli
exposure wall.

\textbf{Part III: The Runtime Architecture (Chapter~\ref{sec-perception}
through Chapter~\ref{sec-placement}).} With a trained policy in hand, we
construct the live, multi-rate software engine that runs aboard the
machine. We follow the continuous flow of signals from photons striking
image sensors, through 3D spatial memory representations and expiring
intent leases, down through receding-horizon trajectory planning and the
bare-metal safety shields that execute on heterogeneous silicon without
starving under bus contention or voltage droop.

\textbf{Part IV: Governance and Defensible Release
(Chapter~\ref{sec-intervention} through Chapter~\ref{sec-frontier}).} In
the final movement, we address the human and societal reality of
physical embodiment. We establish how machines safely hand authority
back to human supervisors without mode confusion, how we break the
assembled system through systematic hardware fault injection across the
qualification ladder, how engineering teams construct auditable
Claim-Argument-Evidence safety cases under emerging international
standards, and how we confront the fundamental epistemic limits of what
AI can and cannot know in the physical world.

\begin{longtable}[]{@{}
  >{\raggedright\arraybackslash}p{(\linewidth - 8\tabcolsep) * \real{0.2000}}
  >{\raggedright\arraybackslash}p{(\linewidth - 8\tabcolsep) * \real{0.2000}}
  >{\raggedright\arraybackslash}p{(\linewidth - 8\tabcolsep) * \real{0.2000}}
  >{\raggedright\arraybackslash}p{(\linewidth - 8\tabcolsep) * \real{0.2000}}
  >{\raggedright\arraybackslash}p{(\linewidth - 8\tabcolsep) * \real{0.2000}}@{}}
\caption{\textbf{The Architecture of Physical AI Systems:
Chapter-by-Chapter Systems Matrix.} Complete structural roadmap of the
textbook across the four foundational parts, mapping each chapter to its
operational focus, primary physical/silicon constraint, and core
engineering mechanism.}\label{tbl-01-book-roadmap}\tabularnewline
\toprule\noalign{}
\begin{minipage}[b]{\linewidth}\raggedright
Part
\end{minipage} & \begin{minipage}[b]{\linewidth}\raggedright
Chapter
\end{minipage} & \begin{minipage}[b]{\linewidth}\raggedright
Subsystem / Focus
\end{minipage} & \begin{minipage}[b]{\linewidth}\raggedright
Core Constraint or Limit
\end{minipage} & \begin{minipage}[b]{\linewidth}\raggedright
Key Systems Mechanism
\end{minipage} \\
\midrule\noalign{}
\endfirsthead
\toprule\noalign{}
\begin{minipage}[b]{\linewidth}\raggedright
Part
\end{minipage} & \begin{minipage}[b]{\linewidth}\raggedright
Chapter
\end{minipage} & \begin{minipage}[b]{\linewidth}\raggedright
Subsystem / Focus
\end{minipage} & \begin{minipage}[b]{\linewidth}\raggedright
Core Constraint or Limit
\end{minipage} & \begin{minipage}[b]{\linewidth}\raggedright
Key Systems Mechanism
\end{minipage} \\
\midrule\noalign{}
\endhead
\bottomrule\noalign{}
\endlastfoot
\textbf{Part I: The Constraints} & Chapter~\ref{sec-boundary} & The
Causal Boundary \& Scope & Idempotency loss; actuation registers &
Proposal--Permission dual-brain contract \\
& Chapter~\ref{sec-body} & Mechanical \& Thermal Dynamics & Reflected
inertia (\(N^2 J\)); Joulean heat & Dynamic stopping envelope
(\(d_{\text{stop}}\)) \\
& Chapter~\ref{sec-brain} & Cognitive Deliberation & Memory wall
(\(M_{\text{weights}}/B_{\text{mem}}\)) & VLA models \& action
chunking \\
& Chapter~\ref{sec-nervous} & Real-Time Computing Bridge & Multi-rate
deadlines \& jitter & Lock-free \texttt{seqlock} \& 5x5 synthesis
grid \\
\textbf{Part II: Teaching the System} & Chapter~\ref{sec-data} &
Endogenous Telemetry & Teleoperation lag; covariate shift & 4-stream
synchronous logging ring buffer \\
& Chapter~\ref{sec-training} & Policy Synthesis Regimes & Compounding
drift (\(\mathcal{O}(T^2 \epsilon)\)) & Diffusion policy \& sim-to-real
domain rand. \\
& Chapter~\ref{sec-evaluation} & Closed-Loop Survival Testing & Butler
\& Finelli exposure wall & Clopper-Pearson exact confidence bounds \\
\textbf{Part III: Running the System} & Chapter~\ref{sec-perception} &
Spatial Perception Waterfall & Sensor transduction lag \& shear &
Geometric clearance insets \& depth cones \\
& Chapter~\ref{sec-memory} & Persistent Spatial Belief & Dynamic
occlusion; transform staleness & TSDF volumetric fusion \& temporal
decay \\
& Chapter~\ref{sec-intent} & Grounded Task Affordances & Semantic
ambiguity \& open horizons & Expiring Intent Leases
(\(\mathcal{L}_{\text{intent}}\)) \\
& Chapter~\ref{sec-planning} & Trajectory Synthesis \& Seams & Seam
torque spikes; execution lag & \(\mathcal{C}^2\) jerk spline blending \&
stop suffixes \\
& Chapter~\ref{sec-enforcement} & Deterministic Safety Shield &
Microcontroller cycle budget (\(1\text{ ms}\)) & Bare-metal 1 kHz CBF-QP
safety filter \\
& Chapter~\ref{sec-placement} & Heterogeneous Silicon QoS & Crossbar
starvation; \(L dI/dt\) droop & AXI QoS bandwidth arbitration \&
fences \\
\textbf{Part IV: Governance \& Assurance} &
Chapter~\ref{sec-intervention} & Shared Human-Machine Autonomy &
Operator reaction delay (\(1\text{ s}\)) & 4-phase handshake \&
\(\mathcal{C}^1\) bumpless transfer \\
& Chapter~\ref{sec-verification} & Hardware Qualification Ladder &
Untested environmental states & SIL \(\to\) PIL \(\to\) HIL fault
injection harness \\
& Chapter~\ref{sec-release} & Safety Cases \& Deployment & Open-world
ODD boundary drift & Claim-Argument-Evidence trees \& canary OTA \\
& Chapter~\ref{sec-frontier} & Epistemic Limits \& Research &
Detectorless physical failure gaps & Multi-timescale governor \&
epistemic bounds \\
\end{longtable}

\section{Systems Stress-Tests}\label{sec-boundary-1-11}

With the core architectural principles, physical archetypes, bedrock
laws, and dual-brain silicon partitioning established, an engineer must
be able to diagnose complex, borderline systems in the field. The
following four systems stress-tests serve as capstone analytical
exercises, challenging you to apply the scope test, physical limiting
budgets, and proposal--permission contracts to real-world edge cases:

\subsection*{Stress-Test 1: The Advisory Observer (Perception Without
Actuation
Authority)}\label{stress-test-1-the-advisory-observer-perception-without-actuation-authority}
\addcontentsline{toc}{subsection}{Stress-Test 1: The Advisory Observer
(Perception Without Actuation Authority)}

\emph{The scenario.} An automotive driver assistance module executes a
50-layer deep convolutional network to detect highway lane markings,
sounding an auditory buzzer whenever the vehicle drifts across a
boundary.

\emph{The systems dilemma.} The system satisfies Criterion 1 (learned
perception model) and Criterion 2 (vehicle motion continuously shifts
camera viewpoints). However, it completely lacks Criterion 3 (delegated
physical authority). Because its only output is an auditory chime, an
algorithmic failure or inference stall cannot directly alter steering
torque or vehicle momentum.

\emph{The systems verdict: Advisory Decision Support.} The system
operates entirely behind glass, relying on the human driver to provide
physical closed-loop actuation.

\subsection*{Stress-Test 2: The Deterministic Autopilot (Actuation
Without Learned
Models)}\label{stress-test-2-the-deterministic-autopilot-actuation-without-learned-models}
\addcontentsline{toc}{subsection}{Stress-Test 2: The Deterministic
Autopilot (Actuation Without Learned Models)}

\emph{The scenario.} An acrobatic quadrotor executes extreme trajectory
tracking through turbulent outdoor wind gusts, calculating motor PWM
commands at \(1000\text{ Hz}\) on a bare-metal microcontroller
(\citeproc{ref-maxwell1868governors}{Maxwell 1868}).

\emph{The systems dilemma.} The system satisfies Criterion 2 (rotor
thrust directly accelerates physical mass) and Criterion 3
(microcontroller latches write directly to motor inverter registers
without human intervention). However, it lacks Criterion 1 (learned
model). The autopilot control law is derived from analytical
Newton-Euler equations and linear state-space feedback whose stability
margins are proven mathematically in closed form.

\emph{The systems verdict: Classical Control Theory.} It requires hard
real-time execution, but does not involve stochastic neural
representations.

\subsection*{Stress-Test 3: The Virtual Agent (Autonomy Without Physical
Feedback)}\label{stress-test-3-the-virtual-agent-autonomy-without-physical-feedback}
\addcontentsline{toc}{subsection}{Stress-Test 3: The Virtual Agent
(Autonomy Without Physical Feedback)}

\emph{The scenario.} An automated high-frequency trading engine executes
thousands of market transactions per second using deep transformer
sequence models to predict price fluctuations.

\emph{The systems dilemma.} The system satisfies Criterion 1 (deep
learned predictive model) and Criterion 3 (automated execution of
transactions without human approval). However, it lacks Criterion 2
(consequential physical feedback). Its outputs alter electronic memory
buffers and order-book records, but do not interact with momentum,
contact friction, or thermodynamic heat.

\emph{The systems verdict: Digital Autonomous Systems.} If an algorithm
diverges, market circuit breakers can freeze accounts, unwind
transactions, and restore clean database states---privileges unavailable
in the physical universe.

\subsection*{Stress-Test 4: The Shared-Autonomy Reflex Gate (The
Neuromuscular Latency
Threshold)}\label{stress-test-4-the-shared-autonomy-reflex-gate-the-neuromuscular-latency-threshold}
\addcontentsline{toc}{subsection}{Stress-Test 4: The Shared-Autonomy
Reflex Gate (The Neuromuscular Latency Threshold)}

\emph{The scenario.} A robotic surgical scalpel applies learned haptic
guidance forces to keep a surgeon's tool tip within a safe anatomical
resection zone.

\emph{The systems dilemma.} The system satisfies Criterion 1 (learned
anatomical spatial models) and Criterion 2 (tool forces directly deform
living tissue). Whether it satisfies Criterion 3 depends strictly on the
actuator force rise time relative to human neuromuscular reflex latency
(\(150\text{--}250\text{ ms}\)). If the haptic actuator can apply forces
within that \(200\text{ ms}\) window, the learned policy exercises
transient, uncontested physical authority over tissue before the human
surgeon can react.

\emph{The systems verdict: Physical AI (Shared Autonomy).} The system
requires a deterministic bare-metal safety filter to clip tool forces
and prevent tissue damage during neural inference anomalies.

\section{Summary}\label{sec-boundary-1-12}

The causal boundary is the physical and electrical interface where
digital information irreversibly converts into thermodynamic action.
Informational abstraction ends at the memory-mapped registers of motor
inverters, valve solenoids, and power stages. Below this boundary,
computational delay is paid in unguided physical displacement, and
released kinetic energy cannot be undone with a software exception
handler.

Surviving this boundary requires an integrated physical container
governed by four bedrock laws:

\begin{enumerate}
\def\labelenumi{\arabic{enumi}.}
\tightlist
\item
  \textbf{The Physical Body} moves mass, exerts contact force, and
  dissipates heat under strict conservation laws.
\item
  \textbf{The Cognitive Brain} executes high-capacity neural models on
  application processors, functioning as an unprivileged proposal
  engine.
\item
  \textbf{The Real-Time Nervous System} executes deterministic safety
  filters on isolated microcontrollers, holding exclusive electrical
  authority over motor registers to verify physical invariants before
  every commit.
\item
  \textbf{The Governance Envelope} enforces decaying operating leases,
  dynamic safety boundaries, and auditable verification cases.
\end{enumerate}

\protect\phantomsection\label{callout-takeawaysux2a-1.10}
\begin{fbxSimple}{callout-takeaways}{Key Chapter Takeaways}{}
\phantomsection\label{callout-takeaways*-1.10}

\begin{itemize}
\tightlist
\item
  \textbf{The Causal Boundary is a Physical Address:} Informational
  abstraction ends at the memory-mapped registers of motor gate drivers.
  Digital bits convert irreversibly into magnetic flux, kinetic
  momentum, and Joule heat (\(I^2 R\)), where software exceptions cannot
  unwind thermodynamic work.
\item
  \textbf{Latency is Paid Directly in Millimeters:} Information decays
  with time (\(t_{\text{age}}\)). Every millisecond of compute and
  transport latency at velocity \(v\) accumulates unguided physical
  displacement (\(\Delta x = \int v \, dt\)), requiring
  sensing-to-actuation budgets to be bounded by strict worst-case tail
  latencies (\(P_{99.9}\)).
\item
  \textbf{The Tri-Archetype Physical Resource Floors:} Physical machines
  divide into Mass-Dominated (Class 1, bound by kinetic stopping
  distance \(v^2/2a_{\max}\)), Contact-Dominated (Class 2, bound by
  contact impedance \(k \Delta x\) and shear stress), and Flow-Dominated
  (Class 3, bound by lumped thermal capacitance and grid phase timing).
\item
  \textbf{The Three-Pillar Scope Test:} An architecture requires the
  full systems discipline of Physical AI if and only if it
  simultaneously combines an unspecifiable learned model, closed-loop
  physical feedback, and delegated actuator authority.
\item
  \textbf{The Four Anatomy Pillars:} Every embodied machine integrates a
  Dynamical Body (Layer 1), a Cognitive Brain (Layer 3), a Real-Time
  Nervous System (Layer 2), and a Governance Envelope (Layer 4).
\item
  \textbf{The Dual-Brain Proposal--Permission Contract:} High-capacity
  neural policies on application processors (Linux MPUs) operate
  strictly as unprivileged proposal engines (\(p_t\)), while isolated
  real-time microcontrollers (bare-metal MCUs) retain exclusive hardware
  authority to evaluate physical barrier invariants before committing
  PWM gate pulses.
\end{itemize}

\end{fbxSimple}

\protect\phantomsection\label{callout-labux2a-1.11}
\begin{fbxSimple}{callout-lab}{Hands-On Lab:}{The Educational Desktop Twin}
\phantomsection\label{callout-lab*-1.11}

To explore these proposal--permission dynamics on an engineer's desk
without requiring industrial factory floors, this course utilizes a
canonical heterogeneous dual-brain System-on-Chip (SoC) as an accessible
desktop reference twin. The architecture pairs an unprivileged Linux
microprocessor (MPU) executing neural vision and planning models with an
isolated, deterministic bare-metal microcontroller (MCU) running
\(1\text{ kHz}\) safety filters over shared SRAM mailboxes. This
desk-bench platform mirrors the exact proposal--permission boundary used
in commercial autonomous vehicles and industrial workcells.

\end{fbxSimple}

The next chapter (Chapter~\ref{sec-body}) transitions from conceptual
boundaries to physical metrology on the hardware floor. By attaching
measuring instruments directly to the compute substrate, the power
electronics, and the mechanical linkages, we derive the governing
equations and measurement disciplines for time, inertia, actuation,
energy, and silicon. That quantitative accounting establishes the firm
numerical budgets that every subsequent chapter must respect.

\part{The Machine Anatomy}

\lettrine{M}{achine} learning was born behind glass. For six decades,
computer science evolved under the protective abstraction of
idempotency: algorithms processed digital symbols in sandboxed memories
where time was elastic, memory was cheap in the cloud, errors printed
stack traces, and computation existed apart from physical consequence.
If an inference pass took an extra hundred milliseconds, a user simply
waited; if a database transaction aborted, the runtime rolled back to a
clean checkpoint; if a model produced an incorrect token, the physical
universe remained untouched.

Physical AI begins at the exact boundary where computation ceases to be
symbolic and becomes physical force. The moment an algorithm writes to a
memory-mapped register on a motor timer, an H-bridge switches, current
surges into copper windings, and mass accelerates through physical space
under momentum and gravity.

At that boundary, all the forgiving conventions of digital software
engineering vanish. Matter will not pause while an operating system
schedules a thread. Inertia does not negotiate with the tail latency of
a neural network. And physical action cannot be undone: kinetic energy
delivered into an environment permanently alters the physical state of
the world, heat accumulates irreversibly in stator windings, and a
mechanical collision cannot be caught by a software exception handler.

To engineer systems that can operate across this divide, Part I
establishes the \textbf{Constraint-Driven Architecture} of physical AI.
Rather than treating an autonomous machine as an abstract software
pipeline, we ground the entire discipline in the fundamental Triad that
comprises every physical machine:

\begin{itemize}
\tightlist
\item
  \textbf{The Body (Chapter 2):} The physical, mechanical, and thermal
  plant. The body obeys the conservation laws of classical mechanics,
  thermodynamics, and electromagnetism rather than software
  instructions. It is governed by five non-negotiable physical budgets:
  sensory measurement freshness, dynamic stopping distance, reflected
  rotor inertia (\(N^2 J\)), stator thermal dissipation (\(I^2 R\)), and
  power delivery voltage droop.
\item
  \textbf{The Brain (Chapter 3):} The learned cognitive deliberative
  engine. We define what the brain is \emph{for}---high-capacity
  open-world semantic generalization, spatial affordance extraction, and
  continuous trajectory synthesis that classical hand-crafted state
  machines cannot achieve. Crucially, we also define what \emph{bounds}
  the brain: the memory wall, Roofline arithmetic intensity limits,
  stochastic fat-tailed latency, and uncalibrated confidence. Because of
  these constraints, the brain must operate as an \textbf{unprivileged
  proposal engine}, emitting advisory action chunks rather than direct
  actuator commands.
\item
  \textbf{The Nervous System (Chapter 4):} The deterministic
  architectural bridge. Sitting between the slow, stochastic brain
  (\(1\text{--}5\text{ Hz}\)) and the fast, unforgiving body
  (\(1000\text{ Hz}\)), the nervous system arbitrates time and
  authority. It translates disparate operational cadences into a
  coherent multi-rate hierarchy, enforces physical invariants at
  microsecond cadences via Control Barrier Functions, and ensures that a
  delay, panic, or hallucination in high-level software never leads to
  unmanaged kinetic damage.
\end{itemize}

Part I provides the mental models, mathematical derivations, and
architectural principles required to understand this Triad. Before we
collect demonstrations, train neural policies, or build runtime
perception stacks, we must first master the physical constraints that
dictate how a machine survives in the material world.



\end{document}
